% Luc Destombe --- licence CC BY-NC-SA 4.0
% https://creativecommons.org/licenses/by-nc-sa/4.0/deed.fr
% Fichier autonome : compiler avec pdflatex, deux fois (signets, nombre de pages).
\documentclass[11pt,a4paper]{article}
\makeatletter
\newif\iflycee@unecolonne
\lycee@unecolonnetrue
\newif\iflycee@cm
\newif\iflycee@palatino

\RequirePackage[utf8]{inputenc}
\RequirePackage[T1]{fontenc}
\RequirePackage[french]{babel}
\RequirePackage{etoolbox}

\newcommand{\ShorthandsOff}{\@ifundefined{shorthandoff}{}{\shorthandoff{;:!?}}}
\newcommand{\ShorthandsOn}{\@ifundefined{shorthandon}{}{\shorthandon{;:!?}}}
\ShorthandsOff
\AtBeginDocument{\ShorthandsOn}
\AtBeginEnvironment{tikzpicture}{\ShorthandsOff}

\RequirePackage{geometry}
\RequirePackage{fancyhdr}
\RequirePackage{lastpage}
\RequirePackage{multicol}
\RequirePackage{array}
\RequirePackage{enumitem}
\RequirePackage{xcolor}
\RequirePackage{needspace}

\RequirePackage{amsmath,amssymb}
\iflycee@cm\else\RequirePackage{fourier}\fi
\DeclareMathAlphabet{\mathcal}{OMS}{cmsy}{m}{n}
\renewcommand{\ttdefault}{lmtt}
\begingroup\catcode`\_=\active \catcode`\^=\active
\global\let_\sb \global\let^\sp
\endgroup
\AtBeginDocument{\catcode`\_=12 \mathcode`\_="8000
  \catcode`\^=12 \mathcode`\^="8000 }
\newlength{\retraitcalculs}
\setlength{\retraitcalculs}{2em}

\RequirePackage{tikz}
\usetikzlibrary{arrows.meta,calc,babel,shapes.misc}
\RequirePackage[most]{tcolorbox}
\tcbuselibrary{skins,breakable}

\definecolor{coulDef}{HTML}{1F4E79}
\definecolor{coulProp}{HTML}{7030A0}
\definecolor{coulRem}{HTML}{C55A11}
\definecolor{coulEx}{HTML}{2E7D32}
\definecolor{coulExo}{HTML}{B03A2E}
\definecolor{coulPart}{HTML}{1F4E79}
\definecolor{coulMeth}{HTML}{1B6E4F}
\definecolor{coulCourbe}{HTML}{0B5ED7}
\definecolor{coulCourbeB}{HTML}{C00000}
\definecolor{coulCourbeC}{HTML}{2E7D32}
\definecolor{coulGrille}{HTML}{8C8C8C}
\definecolor{coulRep}{HTML}{1B6E4F}
\definecolor{coulBar}{HTML}{2E7D32}

\newcommand{\R}{\mathbb{R}}
\newcommand{\Rs}{\mathbb{R}^{*}}
\newcommand{\N}{\mathbb{N}}
\newcommand{\Ns}{\mathbb{N}^{*}}
\newcommand{\Z}{\mathbb{Z}}
\newcommand{\Q}{\mathbb{Q}}
\renewcommand{\le}{\leqslant}
\renewcommand{\ge}{\geqslant}
\renewcommand{\approx}{\simeq}
\newcommand{\vect}[1]{\overrightarrow{#1}}
\newcommand{\norme}[1]{\left\lVert #1 \right\rVert}
\newcommand{\intervalle}[2]{\left[#1\,;#2\right]}
\RequirePackage{cancel}
\renewcommand{\CancelColor}{\color{coulExo}}
\newcommand{\fois}[1]{\times{\color{coulExo}#1}}
\newcommand{\barre}[1]{\cancel{#1}}
\newlength{\lycee@eqh}\newlength{\lycee@eqd}
\newsavebox{\lycee@eqa}\newsavebox{\lycee@eqb}\newsavebox{\lycee@eqc}
\newcommand{\lycee@eqcol}[1]{\hbox{$\renewcommand{\arraystretch}{1.35}%
  \begin{array}[t]{@{}>{\displaystyle}r@{}>{\displaystyle{}}l@{}}#1\end{array}$}}
\newcommand{\lycee@eqmesure}[1]{%
  \ifdim\ht#1>\lycee@eqh\setlength{\lycee@eqh}{\ht#1}\fi
  \ifdim\dp#1>\lycee@eqd\setlength{\lycee@eqd}{\dp#1}\fi}
\newcommand{\lycee@eqtrait}{\hspace{1.2em}\tikz[baseline=0pt]\draw[dash pattern=on 2pt off 2pt,
  line width=0.5pt] (0,\lycee@eqh) -- (0,-\lycee@eqd);\hspace{1.2em}}
\newcommand{\eqproduit}[2]{\par\smallskip\noindent C'est une équation produit :
  $A\times B=0 \iff A=0$ ou $B=0$.\par\nopagebreak\smallskip
  \begingroup\setlength{\lycee@eqh}{0pt}\setlength{\lycee@eqd}{0pt}%
  \setbox\lycee@eqa=\lycee@eqcol{#1}\setbox\lycee@eqb=\lycee@eqcol{#2}%
  \lycee@eqmesure\lycee@eqa\lycee@eqmesure\lycee@eqb%
  \noindent\hspace*{\retraitcalculs}\usebox\lycee@eqa\lycee@eqtrait\usebox\lycee@eqb\endgroup\par\smallskip}
\newcommand{\eqproduittrois}[3]{\par\smallskip\noindent C'est une équation produit :
  $A\times B\times C=0 \iff A=0$ ou $B=0$ ou $C=0$.\par\nopagebreak\smallskip
  \begingroup\setlength{\lycee@eqh}{0pt}\setlength{\lycee@eqd}{0pt}%
  \setbox\lycee@eqa=\lycee@eqcol{#1}\setbox\lycee@eqb=\lycee@eqcol{#2}\setbox\lycee@eqc=\lycee@eqcol{#3}%
  \lycee@eqmesure\lycee@eqa\lycee@eqmesure\lycee@eqb\lycee@eqmesure\lycee@eqc%
  \noindent\hspace*{\retraitcalculs}\usebox\lycee@eqa\lycee@eqtrait\usebox\lycee@eqb\lycee@eqtrait\usebox\lycee@eqc\endgroup
  \par\smallskip}

\tikzset{grille/.style={coulGrille,line width=0.4pt}}
\newcommand{\quadrillage}[4]{\draw[grille,step=1] (#1,#2) grid (#3,#4);}
\tikzset{
  croix/.style={cross out,draw,line width=0.6pt,inner sep=0pt,minimum size=4pt},
  croixdroite/.style={croix,rotate=45,line width=0.8pt},
}
\newcommand{\pt}[3]{\path (#1) node[croix]{} node[#2,font=\small,inner sep=2.5pt]{$#3$};}

\newcommand{\lycee@repere}[4]{%
  \draw[grille,step=1] (#1,#2) grid (#3,#4);
  \draw[-{Stealth[length=2mm]},thick] (#1,0)--({#3+0.4},0) node[below left=-1pt]{$x$};
  \draw[-{Stealth[length=2mm]},thick] (0,#2)--(0,{#4+0.4}) node[below left=-1pt]{$y$};
  \node[below left,font=\scriptsize,inner sep=1.5pt] at (0,0) {$0$};}
\newcommand{\lycee@gradx}[1]{\foreach \k/\l in {#1}{%
  \draw (\k,0.07)--(\k,-0.07) node[below,font=\scriptsize,inner sep=1.5pt]{$\l$};}}
\newcommand{\lycee@grady}[1]{\foreach \k/\l in {#1}{%
  \draw (0.07,\k)--(-0.07,\k) node[left,font=\scriptsize,inner sep=1.5pt]{$\l$};}}
\newcommand{\lycee@intff}[2]{\left[#1\,;#2\right]}
\newcommand{\lycee@intfo}[2]{\left[#1\,;#2\right[}
\newcommand{\lycee@intof}[2]{\left]#1\,;#2\right]}
\newcommand{\lycee@intoo}[2]{\left]#1\,;#2\right[}
\newcommand{\lycee@ens}[1]{\left\{#1\right\}}
\AtBeginDocument{%
  \providecommand{\repere}{\lycee@repere}%
  \providecommand{\gradx}{\lycee@gradx}%
  \providecommand{\grady}{\lycee@grady}%
  \providecommand{\Cf}{\mathcal{C}\sb{f}}%
  \providecommand{\Cg}{\mathcal{C}\sb{g}}%
  \providecommand{\intff}{\lycee@intff}%
  \providecommand{\intfo}{\lycee@intfo}%
  \providecommand{\intof}{\lycee@intof}%
  \providecommand{\intoo}{\lycee@intoo}%
  \providecommand{\ens}{\lycee@ens}}

\newlist{questions}{enumerate}{3}
\setlist[questions,1]{label=\arabic*\textdegree),leftmargin=*,itemsep=2pt,topsep=3pt}
\setlist[questions,2]{label=\alph*),leftmargin=*,itemsep=1pt,topsep=2pt}
\setlist[questions,3]{label=\roman*),leftmargin=*,itemsep=1pt,topsep=1pt}
\setlist[itemize]{label=\textbullet,leftmargin=*,itemsep=2pt,topsep=3pt}

\newcommand{\besoinplace}[1]{\par\penalty\@M\begingroup
  \dimen@=#1\relax\dimen@ii=\pagegoal\advance\dimen@ii-\pagetotal
  \ifdim\dimen@>\dimen@ii\ifdim\dimen@ii>\z@\vfil\fi\break\fi\endgroup}

\newcommand{\lycee@pied}{\thepage/\pageref{LastPage}}
\newcommand{\entete}[2]{%
  \pagestyle{fancy}\fancyhf{}%
  \fancyhead[L]{\footnotesize\color{coulPart}\textbf{#1}}%
  \fancyhead[R]{\footnotesize\color{coulPart}\textbf{#2}}%
  \fancyfoot[C]{\footnotesize\lycee@pied}%
  \renewcommand{\headrulewidth}{0.4pt}%
  \renewcommand{\headrule}{\hbox to\headwidth{%
    \color{coulPart!40}\leaders\hrule height \headrulewidth\hfill}}}

\RequirePackage{ccicons}
\newcommand{\lycee@licence}{{\scriptsize\color{gray}%
  \copyright\ Luc Destombe --- \ccbyncsa\ CC BY-NC-SA 4.0}}
\newif\iflycee@licence \lycee@licencetrue
\newcommand{\sanslicence}{\lycee@licencefalse}
\AtBeginDocument{\iflycee@licence\AddToHookNext{shipout/foreground}{%
  \put(0,0){\rlap{\kern\dimexpr1in+\hoffset+\oddsidemargin+\textwidth\relax
    \llap{\raisebox{-\dimexpr1in+\voffset+\topmargin+\headheight+\headsep
      +\textheight+\footskip\relax}{\lycee@licence}}}}}\fi}

\newcommand{\lycee@site}{faire-des-maths.fr}
\newcommand{\lycee@version}{{\scriptsize\color{gray}\lycee@site
  \ifdefined\lyceeversion\ --- Version~\lyceeversion\fi}}
\AtBeginDocument{\iflycee@licence\AddToHookNext{shipout/foreground}{%
  \put(0,0){\rlap{\kern\dimexpr1in+\hoffset+\oddsidemargin\relax
    \raisebox{-\dimexpr1in+\voffset+\topmargin+\headheight+\headsep
      +\textheight+\footskip\relax}{\lycee@version}}}}\fi}

\newcommand{\formulesserrees}{%
  \setlength{\abovedisplayskip}{3pt}\setlength{\belowdisplayskip}{3pt}%
  \setlength{\abovedisplayshortskip}{2pt}\setlength{\belowdisplayshortskip}{3pt}}

\setlength{\parindent}{0pt}

\RequirePackage[implicit=false,hidelinks,pdfencoding=auto,psdextra,bookmarksopen,
  bookmarksopenlevel=1,pdfcreator={},pdfproducer={}]{hyperref}
\RequirePackage{bookmark}
\hypersetup{pdfauthor={Luc Destombe},pdfkeywords={CC BY-NC-SA 4.0}}
\newcounter{lycee@signet}
\newcommand{\signet}[3][]{\stepcounter{lycee@signet}%
  \edef\lycee@dest{\if\relax\detokenize{#1}\relax
    signet.\arabic{lycee@signet}\else#1\fi}%
  \hypertarget{\lycee@dest}{}\bookmark[level=#2,dest=\lycee@dest]{#3}}
\pdfstringdefDisableCommands{%
  \def\R{R}\def\N{N}\def\Z{Z}\def\Q{Q}\def\vect#1{#1}\def\Cf{Cf}\def\Cg{Cg}\def\fois#1{×#1}%
  \def\barre#1{#1}\def\intervalle#1#2{[#1;#2]}\def\intff#1#2{[#1;#2]}%
  \def\textsuperscript#1{#1}\def\dfrac#1#2{#1/#2}\def\frac#1#2{#1/#2}%
  \def\vec#1{#1⃗}}%
\RequirePackage[official]{eurosym}

\tcbset{exercice corrige/.style={
  enhanced,breakable,before upper=\raggedright,
  colback=white,colframe=coulExo,
  boxrule=0.8pt,arc=2pt,
  left=6pt,right=6pt,top=8pt,bottom=5pt,
  before skip=9pt,after skip=4pt,
  coltitle=white,fonttitle=\bfseries\small,
  attach boxed title to top left={xshift=8pt,yshift=-\tcboxedtitleheight/2},
  boxed title style={colback=coulExo,arc=2pt,boxrule=0pt}}}
\newcounter{exo}
\newtcolorbox{exercice}[1][]{exercice corrige,
  phantom={\signet[exercice-\arabic{exo}]{2}{Exercice \arabic{exo}}},
  title={Exercice \theexo\ifx\relax#1\relax\else{} --- #1\fi}}
\newenvironment{exo}[1][\relax]{\refstepcounter{exo}\begin{exercice}[#1]}{\end{exercice}}

\geometry{top=1.8cm,bottom=1cm,left=1cm,right=1cm,headheight=14pt,footskip=0.6cm}
\renewcommand{\lycee@pied}{\thepage}

\renewcommand{\theexo}{\ifnum\value{exo}<10 0\fi\arabic{exo}}
\tcbset{exercice corrige/.style={
  enhanced,breakable,before upper=\raggedright,
  colback=white,colframe=coulExo,
  boxrule=0.9pt,arc=2pt,
  left=8pt,right=8pt,top=8pt,bottom=6pt,
  before skip=8pt,after skip=6pt,
  coltitle=white,fonttitle=\bfseries,
  attach boxed title to top left={xshift=8pt,yshift=-\tcboxedtitleheight/2},
  boxed title style={colback=coulExo,arc=2pt,boxrule=0pt}}}

\newcommand{\entetefiche}[2]{%
  \begin{center}
    \begin{tcolorbox}[enhanced,width=0.92\linewidth,colback=white,colframe=coulPart,
        boxrule=1.4pt,arc=3pt,halign=center,top=6pt,bottom=6pt]
      {\large\bfseries\color{coulPart} CORRIGÉ --- FICHE D'EXERCICES}\\[3pt]
      {\normalsize\bfseries #1}\\[2pt]
      {\normalsize #2}
    \end{tcolorbox}
  \end{center}}

\newcounter{partie}
\newcommand{\partie}[1]{%
  \stepcounter{partie}%
  \besoinplace{8\baselineskip}\vspace{4pt}%
  \begin{tcolorbox}[enhanced,colback=coulPart!8,colframe=coulPart,boxrule=0pt,
      leftrule=3pt,arc=0pt,left=6pt,right=4pt,top=3pt,bottom=3pt,
      before skip=6pt,after skip=2pt]
    \bfseries\color{coulPart}\Roman{partie} --- #1\signet{1}{\Roman{partie} --- #1}
  \end{tcolorbox}}
\newcommand{\partiesansnum}[1]{%
  \besoinplace{8\baselineskip}\vspace{4pt}%
  \begin{tcolorbox}[enhanced,colback=coulPart!8,colframe=coulPart,boxrule=0pt,
      leftrule=3pt,arc=0pt,left=6pt,right=4pt,top=3pt,bottom=3pt,
      before skip=6pt,after skip=2pt]
    \bfseries\color{coulPart}#1\signet{1}{#1}
  \end{tcolorbox}}

\setlist[questions,1]{label=\arabic*\textdegree),leftmargin=*,itemsep=3pt,topsep=3pt}
\setlist[questions,2]{label=\alph*),leftmargin=*,itemsep=2pt,topsep=2pt}
\setlist[questions,3]{label=\roman*),leftmargin=*,itemsep=1pt,topsep=1pt}
\newlist{expr}{enumerate}{1}
\setlist[expr]{label={\Alph*\ $=$},leftmargin=*,itemsep=2pt,topsep=3pt}

\newcommand{\res}[1]{{\color{coulRep}#1}}
\newcommand{\rem}[1]{\par\smallskip{\small\itshape\color{black!60}$\rightarrow$ #1}}
\newenvironment{calculs}{\par\smallskip\noindent$\displaystyle\begin{aligned}[t]}%
  {\end{aligned}$\par\smallskip}
\newcommand{\justif}[1]{\quad\text{\small\itshape\color{black!60}#1}}
\newcommand{\cas}[1]{\par\smallskip\noindent\textbf{\color{coulExo}#1}\vadjust{\nobreak}}
\newcommand{\explique}[2]{\par\smallskip\noindent
  \begin{minipage}[t]{0.6\linewidth}\raggedright #1\end{minipage}\hfill
  \begin{minipage}[t]{0.37\linewidth}\raggedright\leavevmode
    {\small\itshape\color{black!60}#2}\end{minipage}\par}

\setlength{\parskip}{2pt}

\makeatother
\entete{Chapitre 9 --- Géométrie repérée et produit scalaire --- corrigé}{1\textsuperscript{re} spécialité}
\let\casligne\cas
\renewcommand{\cas}[1]{\casligne{#1}\par\nopagebreak}
\let\partieligne\partie
\renewcommand{\partie}[1]{\besoinplace{16\baselineskip}\partieligne{#1}}
\BeforeBeginEnvironment{exo}{\besoinplace{10\baselineskip}}
\setlength{\parindent}{0pt}

\newcommand{\coord}[2]{\begin{pmatrix}#1\\#2\end{pmatrix}}
\newcommand{\systeme}[2]{\left\{\begin{array}{l}#1\\#2\end{array}\right.}

\begin{document}

\entetefiche{Chapitre 9 --- Géométrie repérée et produit scalaire}{Première spécialité mathématiques}

\partie{Coordonnées d'un vecteur, somme de deux vecteurs}

\begin{exo}[Coordonnées de vecteurs]
\cas{1)}
\explique{$\vec{u}=-2\vec{\imath}+7\vec{\jmath}$, donc $\res{\vec{u}\coord{-2}{7}}$.}{le nombre devant $\vec{\imath}$, puis celui devant $\vec{\jmath}$}
\explique{$\vec{v}=0\,\vec{\imath}+5\vec{\jmath}$, donc $\res{\vec{v}\coord{0}{5}}$.}{pas de $\vec{\imath}$ : l'abscisse est nulle}
\cas{2) a)}
\begin{calculs}
\vect{AB} & \coord{x_{B}-x_{A}}{y_{B}-y_{A}} &&\justif{arrivée moins départ}\\
\vect{AB} & \coord{6-2}{1-5} &&\justif{on remplace par les coordonnées}\\
\vect{AB} & \res{\coord{4}{-4}} &&\\
\vect{BA} & \res{\coord{-4}{4}} &&\justif{l'opposé de $\vect{AB}$ : le trajet inverse}
\end{calculs}
\cas{2) b)}
\begin{calculs}
\vect{AB} & \coord{1-(-3)}{-4-2} &&\justif{arrivée moins départ}\\
\vect{AB} & \res{\coord{4}{-6}} &&\\
\vect{BA} & \res{\coord{-4}{6}} &&\justif{l'opposé de $\vect{AB}$}
\end{calculs}
\cas{2) c)}
\begin{calculs}
\vect{AB} & \coord{-7-(-1)}{-2-(-5)} &&\justif{attention aux doubles signes « moins »}\\
\vect{AB} & \res{\coord{-6}{3}} &&\\
\vect{BA} & \res{\coord{6}{-3}} &&\justif{l'opposé de $\vect{AB}$}
\end{calculs}
\cas{2) d)}
\begin{calculs}
x_{B}-x_{A} &= -1-\dfrac{1}{3} &&\justif{l'abscisse de $\vect{AB}$}\\
x_{B}-x_{A} &= -\dfrac{1\fois{3}}{1\fois{3}}-\dfrac{1}{3} &&\justif{même dénominateur}\\
x_{B}-x_{A} &= -\dfrac{4}{3} &&\\
y_{B}-y_{A} &= \dfrac{5}{3}-2 &&\justif{l'ordonnée de $\vect{AB}$}\\
y_{B}-y_{A} &= \dfrac{5}{3}-\dfrac{2\fois{3}}{1\fois{3}} &&\\
y_{B}-y_{A} &= -\dfrac{1}{3} &&
\end{calculs}
\explique{Donc $\res{\vect{AB}\coord{-\frac{4}{3}}{-\frac{1}{3}}}$ et $\res{\vect{BA}\coord{\frac{4}{3}}{\frac{1}{3}}}$.}{l'opposé : on change les deux signes}
\end{exo}

\begin{exo}[Calculer avec les coordonnées]
\cas{1)}
\begin{calculs}
\vec{u}+\vec{v} & \coord{2+(-3)}{-5+1} &&\justif{on ajoute les abscisses, puis les ordonnées}\\
\vec{u}+\vec{v} & \res{\coord{-1}{-4}} &&\\
\vec{u}-\vec{v} & \coord{2-(-3)}{-5-1} &&\justif{on soustrait, dans l'ordre}\\
\vec{u}-\vec{v} & \res{\coord{5}{-6}} &&\\
\vec{v}-\vec{u} & \res{\coord{-5}{6}} &&\justif{l'opposé de $\vec{u}-\vec{v}$}
\end{calculs}
\cas{2)}
\explique{$\vec{u}+\vec{t}=\vec{w}$ équivaut à $\vec{t}=\vec{w}-\vec{u}$.}{on retranche $\vec{u}$ des deux côtés, comme dans une équation}
\begin{calculs}
\vec{t} & \coord{4-2}{0-(-5)} &&\justif{coordonnées de $\vec{w}-\vec{u}$}\\
\vec{t} & \res{\coord{2}{5}} &&
\end{calculs}
\rem{On vérifie : $\vec{u}+\vec{t}$ a pour coordonnées $(2+2\,;-5+5)$, soit $(4\,;0)$, celles de $\vec{w}$.}
\end{exo}

\begin{exo}[Relation de Chasles]
\cas{1)}
\begin{calculs}
\vect{AB} & \coord{4-0}{1-3} &&\justif{arrivée moins départ}\\
\vect{AB} & \coord{4}{-2} &&\\
\vect{BC} & \coord{-2-4}{-1-1} &&\\
\vect{BC} & \coord{-6}{-2} &&\\
\vect{AB}+\vect{BC} & \res{\coord{-2}{-4}} &&\justif{$4-6$ et $-2-2$}\\
\vect{AC} & \coord{-2-0}{-1-3} &&\\
\vect{AC} & \res{\coord{-2}{-4}} &&
\end{calculs}
\explique{On constate que $\res{\vect{AB}+\vect{BC}=\vect{AC}}$ : c'est la relation de Chasles.}{deux vecteurs qui ont les mêmes coordonnées sont égaux}
\cas{2)}
\begin{calculs}
\vect{AB}-\vect{AC} & \coord{4-(-2)}{-2-(-4)} &&\\
\vect{AB}-\vect{AC} & \coord{6}{2} &&\\
\vect{CB} & \coord{4-(-2)}{1-(-1)} &&\justif{on essaie $\vect{CB}$}\\
\vect{CB} & \coord{6}{2} &&
\end{calculs}
\explique{On retrouve $\res{\vect{AB}-\vect{AC}=\vect{CB}}$.}{en effet, $\vect{AB}-\vect{AC}=\vect{CA}+\vect{AB}$, soit $\vect{CB}$ par la relation de Chasles}
\end{exo}

\begin{exo}[Trajet d'un drone]
\cas{1)}
\explique{Le point d'arrivée d'un déplacement a pour coordonnées celles du point de départ, augmentées de celles du vecteur.}{$\vect{OM_{1}}=\vec{u}$, puis $\vect{M_{1}M_{2}}=\vec{v}$\dots}
\begin{calculs}
M_{1} & (0+3\,;0+2) &&\justif{on part de $O(0\,;0)$ selon $\vec{u}$}\\
M_{1} & \res{(3\,;2)} &&\\
M_{2} & (3-1\,;2+4) &&\justif{selon $\vec{v}$}\\
M_{2} & \res{(2\,;6)} &&\\
M_{3} & (2+2\,;6-3) &&\justif{selon $\vec{w}$}\\
M_{3} & \res{(4\,;3)} &&
\end{calculs}
\cas{2)}
\begin{calculs}
\vec{u}+\vec{v}+\vec{w} & \coord{3-1+2}{2+4-3} &&\\
\vec{u}+\vec{v}+\vec{w} & \res{\coord{4}{3}} &&
\end{calculs}
\explique{C'est le vecteur $\vect{OM_{3}}$ : il représente \res{le déplacement total}, du départ au point d'arrivée.}{relation de Chasles : $\vect{OM_{1}}+\vect{M_{1}M_{2}}+\vect{M_{2}M_{3}}=\vect{OM_{3}}$}
\cas{3)}
\explique{Le déplacement de vecteur $\vect{M_{3}O}$, soit $\res{\coord{-4}{-3}}$, ramène le drone à son point de départ.}{l'opposé du déplacement total}
\end{exo}

\partie{Parallélogramme et milieu d'un segment}

\begin{exo}[Point défini par une égalité de vecteurs]
\cas{a)}
\begin{calculs}
\vect{AE} & \coord{x-1}{y+2} &&\justif{on pose $E(x\,;y)$ ; $y-(-2)$}\\
x-1 &= -3 &&\justif{mêmes abscisses que $\vec{u}$}\\
x &= -2 &&\\
y+2 &= 4 &&\justif{mêmes ordonnées}\\
y &= 2 &&
\end{calculs}
\explique{Donc $\res{E(-2\,;2)}$.}{}
\cas{b)}
\explique{$ABCD$ est un parallélogramme quand $\vect{AB}=\vect{DC}$.}{propriété 3 : attention à l'ordre des lettres}
\begin{calculs}
\vect{AB} & \coord{3}{3} &&\justif{$4-1$ et $1-(-2)$}\\
\vect{DC} & \coord{2-x}{5-y} &&\justif{on pose $D(x\,;y)$}\\
2-x &= 3 &&\\
x &= -1 &&\\
5-y &= 3 &&\\
y &= 2 &&
\end{calculs}
\explique{Donc $\res{D(-1\,;2)}$.}{}
\cas{c)}
\begin{calculs}
\vect{BC} & \coord{-2}{4} &&\justif{$2-4$ et $5-1$}\\
\vect{FA} & \coord{1-x}{-2-y} &&\justif{on pose $F(x\,;y)$}\\
1-x &= -2 &&\\
x &= 3 &&\\
-2-y &= 4 &&\\
y &= -6 &&
\end{calculs}
\explique{Donc $\res{F(3\,;-6)}$.}{}
\end{exo}

\begin{exo}[Milieux]
\cas{1)}
\begin{calculs}
I & \left(\dfrac{x_{R}+x_{S}}{2}\,;\dfrac{y_{R}+y_{S}}{2}\right) &&\justif{la moyenne des abscisses, puis des ordonnées}\\
I & \left(\dfrac{-4+2}{2}\,;\dfrac{3+7}{2}\right) &&\\
I & \res{(-1\,;5)} &&
\end{calculs}
\cas{2)}
\explique{On pose $T(x\,;y)$ et on écrit que $K$ est le milieu de $[RT]$.}{les inconnues sont les coordonnées de $T$}
\begin{calculs}
\dfrac{-4+x}{2} &= 1 &&\justif{abscisse du milieu}\\
-4+x &= 2 &&\justif{on multiplie par $2$}\\
x &= 6 &&\\
\dfrac{3+y}{2} &= -2 &&\justif{ordonnée du milieu}\\
3+y &= -4 &&\\
y &= -7 &&
\end{calculs}
\explique{Donc $\res{T(6\,;-7)}$.}{}
\cas{3)}
\begin{calculs}
x_{E}+x_{F} &= \dfrac{1}{3}+\dfrac{1\fois{3}}{1\fois{3}} &&\justif{même dénominateur}\\
x_{E}+x_{F} &= \dfrac{4}{3} &&\\
\dfrac{x_{E}+x_{F}}{2} &= \dfrac{4}{3}\times\dfrac{1}{2} &&\justif{diviser par $2$, c'est multiplier par $\frac{1}{2}$}\\
\dfrac{x_{E}+x_{F}}{2} &= \dfrac{2}{3} &&\\
\dfrac{y_{E}+y_{F}}{2} &= \dfrac{2-1}{2} &&\\
\dfrac{y_{E}+y_{F}}{2} &= \dfrac{1}{2} &&
\end{calculs}
\explique{Le milieu de $[EF]$ est le point $\res{\left(\frac{2}{3}\,;\frac{1}{2}\right)}$.}{}
\end{exo}

\begin{exo}[Parallélogramme ou non ?]
\cas{a)}
\explique{Les diagonales de $KLMN$ sont $[KM]$ et $[LN]$.}{on relie les sommets non consécutifs}
\begin{calculs}
\text{milieu de }[KM] & \left(\dfrac{-2+5}{2}\,;\dfrac{-1+4}{2}\right) &&\\
\text{milieu de }[KM] & \left(\dfrac{3}{2}\,;\dfrac{3}{2}\right) &&\\
\text{milieu de }[LN] & \left(\dfrac{3+0}{2}\,;\dfrac{0+3}{2}\right) &&\\
\text{milieu de }[LN] & \left(\dfrac{3}{2}\,;\dfrac{3}{2}\right) &&
\end{calculs}
\explique{Les diagonales ont le même milieu : $KLMN$ \res{est un parallélogramme}.}{un quadrilatère dont les diagonales ont le même milieu est un parallélogramme}
\cas{b)}
\begin{calculs}
\text{milieu de }[LV] & \left(\dfrac{3+1}{2}\,;\dfrac{0+3}{2}\right) &&\justif{les diagonales sont $[KM]$ et $[LV]$}\\
\text{milieu de }[LV] & \left(2\,;\dfrac{3}{2}\right) &&
\end{calculs}
\explique{Ce milieu n'est pas celui de $[KM]$ : $KLMV$ \res{n'est pas un parallélogramme}.}{les abscisses $2$ et $\frac{3}{2}$ diffèrent}
\end{exo}

\begin{exo}[Quatrième sommet]
\cas{1)}
\explique{$ABCD$ est un parallélogramme quand $\vect{AB}=\vect{DC}$ ; on pose $D(x\,;y)$.}{propriété 3}
\begin{calculs}
\vect{AB} & \coord{4}{-2} &&\justif{$3-(-1)$ et $-1-1$}\\
\vect{DC} & \coord{4-x}{2-y} &&\\
4-x &= 4 &&\\
x &= 0 &&\\
2-y &= -2 &&\\
y &= 4 &&
\end{calculs}
\explique{Donc $\res{D(0\,;4)}$.}{}
\cas{2)}
\begin{calculs}
O & \left(\dfrac{-1+4}{2}\,;\dfrac{1+2}{2}\right) &&\justif{le centre est le milieu de la diagonale $[AC]$}\\
O & \res{\left(\dfrac{3}{2}\,;\dfrac{3}{2}\right)} &&
\end{calculs}
\rem{On peut vérifier avec l'autre diagonale : le milieu de $[BD]$ est $\left(\frac{3+0}{2}\,;\frac{-1+4}{2}\right)$, le même point.}
\cas{3)}
\begin{calculs}
\vect{AD} & \coord{1}{3} &&\justif{$0-(-1)$ et $4-1$}\\
\vect{AB}+\vect{AD} & \coord{4+1}{-2+3} &&\\
\vect{AB}+\vect{AD} & \res{\coord{5}{1}} &&\\
\vect{AC} & \coord{4-(-1)}{2-1} &&\\
\vect{AC} & \res{\coord{5}{1}} &&
\end{calculs}
\explique{On retrouve la \res{règle du parallélogramme} : $\vect{AB}+\vect{AD}=\vect{AC}$.}{la somme de deux côtés issus de $A$ est la diagonale issue de $A$}
\cas{4)}
\explique{On pose $E(x\,;y)$ ; $B$ est le milieu de $[AE]$.}{}
\begin{calculs}
\dfrac{-1+x}{2} &= 3 &&\justif{abscisse du milieu}\\
x &= 7 &&\justif{$-1+x=6$}\\
\dfrac{1+y}{2} &= -1 &&\justif{ordonnée du milieu}\\
y &= -3 &&\justif{$1+y=-2$}
\end{calculs}
\explique{Donc $\res{E(7\,;-3)}$.}{on peut aussi écrire $\vect{BE}=\vect{AB}$}
\end{exo}

\partie{Produit d'un vecteur par un nombre réel, colinéarité}

\begin{exo}[Vecteurs colinéaires ou non]
\cas{1) a)}
\begin{calculs}
\det(\vec{u}\,;\vec{v}) &= 3\times 12-(-4)\times(-9) &&\justif{$xy'-yx'$}\\
\det(\vec{u}\,;\vec{v}) &= 36-36 &&\\
\det(\vec{u}\,;\vec{v}) &= 0 &&
\end{calculs}
\explique{$\vec{u}$ et $\vec{v}$ \res{sont colinéaires}.}{d'ailleurs $\vec{v}=-3\vec{u}$}
\cas{1) b)}
\begin{calculs}
\det(\vec{u}\,;\vec{w}) &= 5\times 3-2\times 7 &&\justif{$xy'-yx'$}\\
\det(\vec{u}\,;\vec{w}) &= 1 &&
\end{calculs}
\explique{Le déterminant n'est pas nul : $\vec{u}$ et $\vec{w}$ \res{ne sont pas colinéaires}.}{}
\cas{1) c)}
\begin{calculs}
\det(\vec{a}\,;\vec{b}) &= \sqrt{2}\times\sqrt{2}-2\times 1 &&\justif{$xy'-yx'$}\\
\det(\vec{a}\,;\vec{b}) &= 2-2 &&\justif{$\sqrt{2}\times\sqrt{2}=2$}\\
\det(\vec{a}\,;\vec{b}) &= 0 &&
\end{calculs}
\explique{$\vec{a}$ et $\vec{b}$ \res{sont colinéaires}.}{d'ailleurs $\vec{a}=\sqrt{2}\,\vec{b}$}
\cas{2)}
\begin{calculs}
\det(\vec{r}\,;\vec{s}) &= m\times m-3\times 27 &&\justif{$xy'-yx'$}\\
\det(\vec{r}\,;\vec{s}) &= m^{2}-81 &&
\end{calculs}
\explique{$\vec{r}$ et $\vec{s}$ sont colinéaires quand $m^{2}-81=0$, soit $m^{2}=81$.}{le déterminant doit être nul}
\explique{Donc $\res{m=9}$ ou $\res{m=-9}$.}{deux nombres ont pour carré $81$ : ne pas oublier $-9$}
\end{exo}

\begin{exo}[Combinaisons de vecteurs]
\cas{1)}
\begin{calculs}
3\vec{u} & \coord{6}{-3} &&\justif{on multiplie chaque coordonnée par $3$}\\
2\vec{v} & \coord{-6}{8} &&\\
3\vec{u}-2\vec{v} & \coord{6-(-6)}{-3-8} &&\\
3\vec{u}-2\vec{v} & \res{\coord{12}{-11}} &&\\
\frac{1}{2}\vec{u} & \coord{1}{-\frac{1}{2}} &&\justif{la moitié de chaque coordonnée}\\
\frac{1}{2}\vec{u}+\vec{v} & \coord{1-3}{-\frac{1}{2}+4} &&\\
\frac{1}{2}\vec{u}+\vec{v} & \res{\coord{-2}{\frac{7}{2}}} &&\justif{$-\frac{1}{2}+\frac{8}{2}$}
\end{calculs}
\cas{2)}
\explique{$2\vec{w}+\vec{u}=\vec{v}$ équivaut à $2\vec{w}=\vec{v}-\vec{u}$, puis à $\vec{w}=\frac{1}{2}\left(\vec{v}-\vec{u}\right)$.}{on isole $\vec{w}$ comme dans une équation}
\begin{calculs}
\vec{v}-\vec{u} & \coord{-3-2}{4-(-1)} &&\\
\vec{v}-\vec{u} & \coord{-5}{5} &&\\
\vec{w} & \res{\coord{-\frac{5}{2}}{\frac{5}{2}}} &&\justif{la moitié de chaque coordonnée}
\end{calculs}
\end{exo}

\begin{exo}[Parallélisme et alignement]
\cas{1)}
\begin{calculs}
\vect{AB} & \coord{4}{2} &&\justif{$2-(-2)$ et $3-1$}\\
\vect{CD} & \coord{8}{4} &&\justif{$8-0$ et $2-(-2)$}\\
\det(\vect{AB}\,;\vect{CD}) &= 4\times 4-2\times 8 &&\justif{$xy'-yx'$}\\
\det(\vect{AB}\,;\vect{CD}) &= 0 &&
\end{calculs}
\explique{Les vecteurs sont colinéaires : les droites $(AB)$ et $(CD)$ \res{sont parallèles}.}{propriété 5}
\cas{2)}
\begin{calculs}
\vect{AE} & \coord{10}{5} &&\justif{$8-(-2)$ et $6-1$}\\
\det(\vect{AB}\,;\vect{AE}) &= 4\times 5-2\times 10 &&\\
\det(\vect{AB}\,;\vect{AE}) &= 0 &&
\end{calculs}
\explique{$\vect{AB}$ et $\vect{AE}$ sont colinéaires : les points $A$, $B$ et $E$ \res{sont alignés}.}{deux vecteurs de même origine $A$}
\begin{calculs}
\vect{AC} & \coord{2}{-3} &&\justif{$0-(-2)$ et $-2-1$}\\
\det(\vect{AB}\,;\vect{AC}) &= 4\times(-3)-2\times 2 &&\\
\det(\vect{AB}\,;\vect{AC}) &= -16 &&
\end{calculs}
\explique{Le déterminant n'est pas nul : $A$, $B$ et $C$ \res{ne sont pas alignés}.}{}
\cas{3)}
\explique{$M$ est sur $(AB)$ quand $\vect{AB}$ et $\vect{AM}$ sont colinéaires.}{même méthode, avec une inconnue}
\begin{calculs}
\vect{AM} & \coord{12}{y-1} &&\justif{$10-(-2)$}\\
\det(\vect{AB}\,;\vect{AM}) &= 4(y-1)-2\times 12 &&\\
\det(\vect{AB}\,;\vect{AM}) &= 4y-28 &&\justif{$4y-4-24$}\\
4y-28 &= 0 &&\justif{le déterminant doit être nul}\\
y &= 7 &&
\end{calculs}
\explique{Le point $M$ a pour ordonnée $\res{y=7}$.}{}
\end{exo}

\begin{exo}[Droite des milieux]
\cas{1)}
\begin{calculs}
I & \left(\dfrac{-2+2}{2}\,;\dfrac{0+6}{2}\right) &&\justif{milieu de $[AC]$}\\
I & \res{(0\,;3)} &&\\
J & \left(\dfrac{4+2}{2}\,;\dfrac{2+6}{2}\right) &&\justif{milieu de $[BC]$}\\
J & \res{(3\,;4)} &&
\end{calculs}
\cas{2)}
\begin{calculs}
\vect{IJ} & \coord{3}{1} &&\justif{$3-0$ et $4-3$}\\
\vect{AB} & \coord{6}{2} &&\justif{$4-(-2)$ et $2-0$}\\
\frac{1}{2}\vect{AB} & \coord{3}{1} &&\justif{la moitié de chaque coordonnée}
\end{calculs}
\explique{Les deux vecteurs ont les mêmes coordonnées, donc $\res{\vect{IJ}=\frac{1}{2}\vect{AB}}$.}{}
\cas{3)}
\explique{$\vect{IJ}$ et $\vect{AB}$ sont colinéaires, donc les droites $(IJ)$ et $(AB)$ \res{sont parallèles}.}{propriété 5 ; on retrouve le théorème des milieux}
\cas{4)}
\begin{calculs}
\vect{AE} & \coord{-3}{-1} &&\justif{$-5-(-2)$ et $-1-0$}\\
\det(\vect{AB}\,;\vect{AE}) &= 6\times(-1)-2\times(-3) &&\\
\det(\vect{AB}\,;\vect{AE}) &= 0 &&\justif{$-6+6$}
\end{calculs}
\explique{$A$, $B$ et $E$ sont alignés : $E$ \res{est sur la droite $(AB)$}.}{}
\end{exo}

\partie{Norme d'un vecteur, équation d'un cercle}

\begin{exo}[Normes et distances]
\cas{1)}
\begin{calculs}
\norme{\vec{u}} &= \sqrt{x^{2}+y^{2}} &&\justif{repère orthonormé}\\
\norme{\vec{u}} &= \sqrt{(-7)^{2}+24^{2}} &&\\
\norme{\vec{u}} &= \sqrt{625} &&\justif{$49+576$}\\
\norme{\vec{u}} &= \res{25} &&
\end{calculs}
\cas{2)}
\begin{calculs}
\vect{CD} & \coord{-4}{8} &&\justif{$-3-1$ et $6-(-2)$}\\
CD &= \sqrt{(-4)^{2}+8^{2}} &&\\
CD &= \sqrt{80} &&\justif{$16+64$}\\
CD &= \res{4\sqrt{5}} &&\justif{$\sqrt{16}\times\sqrt{5}$ : on cherche le plus grand carré qui divise $80$}
\end{calculs}
\cas{3) a)}
\explique{On compare les carrés des trois côtés.}{pas besoin des racines carrées pour comparer}
\begin{calculs}
EF^{2} &= 2^{2}+4^{2} &&\justif{$\vect{EF}\coord{2}{4}$}\\
EF^{2} &= 20 &&\\
FG^{2} &= (-4)^{2}+2^{2} &&\justif{$\vect{FG}\coord{-4}{2}$}\\
FG^{2} &= 20 &&\\
EG^{2} &= (-2)^{2}+6^{2} &&\justif{$\vect{EG}\coord{-2}{6}$}\\
EG^{2} &= 40 &&
\end{calculs}
\explique{$EF=FG$ et $EF^{2}+FG^{2}=EG^{2}$ : le triangle $EFG$ est \res{rectangle et isocèle en $F$}.}{réciproque du théorème de Pythagore}
\cas{3) b)}
\begin{calculs}
KL^{2} &= 4^{2}+3^{2} &&\justif{$\vect{KL}\coord{4}{3}$}\\
KL^{2} &= 25 &&\\
KM^{2} &= (-3)^{2}+(-4)^{2} &&\justif{$\vect{KM}\coord{-3}{-4}$}\\
KM^{2} &= 25 &&\\
LM^{2} &= (-7)^{2}+(-7)^{2} &&\justif{$\vect{LM}\coord{-7}{-7}$}\\
LM^{2} &= 98 &&
\end{calculs}
\explique{$KL=KM$ : le triangle $KLM$ est \res{isocèle en $K$} ; il n'est pas rectangle, car $25+25$ ne vaut pas $98$.}{le plus grand côté est $[LM]$ : c'est lui qu'on teste}
\end{exo}

\begin{exo}[Équations de cercles]
\cas{1) a)}
\explique{$(x-4)^{2}+(y-(-1))^{2}=3^{2}$ : centre $\res{(4\,;-1)}$, rayon $\res{3}$.}{forme $(x-a)^{2}+(y-b)^{2}=r^{2}$ : $y+1$, c'est $y-(-1)$}
\cas{1) b)}
\explique{$(x-0)^{2}+(y-2)^{2}=\big(\sqrt{5}\big)^{2}$ : centre $\res{(0\,;2)}$, rayon $\res{\sqrt{5}}$.}{le rayon est la racine carrée du second membre}
\cas{2)}
\explique{$\mathcal{C}_{1}$ : $(x-(-3))^{2}+(y-2)^{2}=5^{2}$, soit $\res{(x+3)^{2}+(y-2)^{2}=25}$.}{théorème 2}
\cas{3)}
\begin{calculs}
AB^{2} &= (-1-2)^{2}+(5-1)^{2} &&\justif{le rayon est $AB$ ; son carré suffit}\\
AB^{2} &= 25 &&\justif{$9+16$}
\end{calculs}
\explique{$\mathcal{C}_{2}$ : $\res{(x-2)^{2}+(y-1)^{2}=25}$.}{}
\begin{calculs}
(5-2)^{2}+(5-1)^{2} &= 9+16 &&\justif{on remplace par les coordonnées de $C$}\\
(5-2)^{2}+(5-1)^{2} &= 25 &&\\
(-2-2)^{2}+(3-1)^{2} &= 16+4 &&\justif{coordonnées de $D$}\\
(-2-2)^{2}+(3-1)^{2} &= 20 &&
\end{calculs}
\explique{$C$ \res{est sur $\mathcal{C}_{2}$} ; $D$ \res{n'y est pas} ($20$ n'est pas $25$).}{}
\cas{4)}
\explique{Le centre $I$ est le milieu de $[EF]$, le rayon est la moitié de $EF$.}{}
\begin{calculs}
I & \left(\dfrac{-4+2}{2}\,;\dfrac{3-5}{2}\right) &&\\
I & (-1\,;-1) &&\\
IE^{2} &= (-4+1)^{2}+(3+1)^{2} &&\justif{le rayon est $IE$}\\
IE^{2} &= 25 &&\justif{$9+16$}
\end{calculs}
\explique{$\mathcal{C}_{3}$ : $\res{(x+1)^{2}+(y+1)^{2}=25}$.}{}
\end{exo}

\begin{exo}[Reconnaître un cercle]
\cas{a)}
\explique{On complète les carrés : $x^{2}-6x=(x-3)^{2}-9$ et $y^{2}+4y=(y+2)^{2}-4$.}{forme canonique du chapitre 3 : $x^{2}-6x$ est le début de $(x-3)^{2}$}
\begin{calculs}
x^{2}+y^{2}-6x+4y-3 &= 0 &&\justif{on recopie l'équation}\\
(x-3)^{2}-9+(y+2)^{2}-4-3 &= 0 &&\justif{on remplace}\\
(x-3)^{2}+(y+2)^{2} &= 16 &&\justif{$9+4+3$}
\end{calculs}
\explique{C'est le cercle de centre $\res{(3\,;-2)}$ et de rayon $\res{4}$.}{$16$ est positif}
\cas{b)}
\begin{calculs}
x^{2}+y^{2}+8x-2y+8 &= 0 &&\\
(x+4)^{2}-16+(y-1)^{2}-1+8 &= 0 &&\justif{$x^{2}+8x=(x+4)^{2}-16$ et $y^{2}-2y=(y-1)^{2}-1$}\\
(x+4)^{2}+(y-1)^{2} &= 9 &&\justif{$16+1-8$}
\end{calculs}
\explique{C'est le cercle de centre $\res{(-4\,;1)}$ et de rayon $\res{3}$.}{}
\cas{c)}
\begin{calculs}
x^{2}+y^{2}-2x+6y+11 &= 0 &&\\
(x-1)^{2}-1+(y+3)^{2}-9+11 &= 0 &&\\
(x-1)^{2}+(y+3)^{2} &= -1 &&\justif{$1+9-11$}
\end{calculs}
\explique{Une somme de carrés ne peut pas valoir $-1$ : \res{ce n'est pas un cercle} (aucun point ne vérifie l'équation).}{un second membre négatif n'est pas un rayon au carré}
\cas{Milieu}
\begin{calculs}
\text{milieu} & \left(\dfrac{3+(-4)}{2}\,;\dfrac{-2+1}{2}\right) &&\justif{centres $(3\,;-2)$ et $(-4\,;1)$}\\
\text{milieu} & \res{\left(-\dfrac{1}{2}\,;-\dfrac{1}{2}\right)} &&
\end{calculs}
\end{exo}

\begin{exo}[Zone couverte par une antenne]
\cas{1)}
\explique{$\mathcal{C}$ : $\res{(x-2)^{2}+(y-3)^{2}=25}$.}{centre $\Omega(2\,;3)$, rayon $5$}
\cas{2)}
\begin{calculs}
\Omega V^{2} &= (5-2)^{2}+(7-3)^{2} &&\justif{carré de la distance à l'antenne}\\
\Omega V^{2} &= 25 &&\justif{$9+16$}\\
\Omega W^{2} &= (-2-2)^{2}+(-1-3)^{2} &&\\
\Omega W^{2} &= 32 &&\justif{$16+16$}
\end{calculs}
\explique{$\Omega V=5$ : le village $V$ est sur le bord de la zone, il \res{est couvert} ; $\Omega W>5$ : $W$ \res{n'est pas couvert}.}{un point est couvert quand sa distance à $\Omega$ est au plus $5$, soit son carré au plus $25$}
\cas{3)}
\begin{calculs}
\dfrac{5+x}{2} &= 2 &&\justif{$V'(x\,;y)$, $\Omega$ milieu de $[VV']$}\\
x &= -1 &&\justif{$5+x=4$}\\
\dfrac{7+y}{2} &= 3 &&\\
y &= -1 &&\justif{$7+y=6$}
\end{calculs}
\explique{Donc $\res{V'(-1\,;-1)}$.}{$V'$ est diamétralement opposé à $V$ : il est aussi sur $\mathcal{C}$}
\cas{4)}
\begin{calculs}
(x-2)^{2}+(-1-3)^{2} &= 25 &&\justif{on remplace $y$ par $-1$}\\
(x-2)^{2}+16 &= 25 &&\\
(x-2)^{2} &= 9 &&
\end{calculs}
\explique{$x-2=3$ ou $x-2=-3$, soit $x=5$ ou $x=-1$.}{deux nombres ont pour carré $9$}
\explique{Les points de la route sur $\mathcal{C}$ sont $\res{(5\,;-1)}$ et $\res{(-1\,;-1)}$ ; le second est \res{le point $V'$}.}{}
\end{exo}

\partie{Définition du produit scalaire}

\begin{exo}[Produit scalaire avec un angle]
\cas{a)}
\begin{calculs}
\vec{u}\cdot\vec{v} &= \norme{\vec{u}}\times\norme{\vec{v}}\times\cos(\alpha) &&\justif{définition 4}\\
\vec{u}\cdot\vec{v} &= 4\times 5\times\cos(60^{\circ}) &&\\
\vec{u}\cdot\vec{v} &= 20\times\dfrac{1}{2} &&\\
\vec{u}\cdot\vec{v} &= \res{10} &&
\end{calculs}
\cas{b)}
\begin{calculs}
\vect{AB}\cdot\vect{AC} &= AB\times AC\times\cos\big(\widehat{BAC}\big) &&\justif{deux vecteurs issus de $A$}\\
\vect{AB}\cdot\vect{AC} &= 6\times 2\times\cos(150^{\circ}) &&\\
\vect{AB}\cdot\vect{AC} &= 12\times\left(-\dfrac{\sqrt{3}}{2}\right) &&\justif{$\cos(150^{\circ})=-\cos(30^{\circ})$}\\
\vect{AB}\cdot\vect{AC} &= \res{-6\sqrt{3}} &&\justif{négatif : l'angle est obtus}
\end{calculs}
\cas{c)}
\begin{calculs}
\vect{DE}\cdot\vect{DF} &= 3\times 4\sqrt{2}\times\cos(45^{\circ}) &&\\
\vect{DE}\cdot\vect{DF} &= 12\sqrt{2}\times\dfrac{\sqrt{2}}{2} &&\\
\vect{DE}\cdot\vect{DF} &= \res{12} &&\justif{$\sqrt{2}\times\sqrt{2}=2$}
\end{calculs}
\cas{d)}
\begin{calculs}
\vect{KL}\cdot\vect{KM} &= 8\times 8\times\cos(60^{\circ}) &&\justif{triangle équilatéral : angles de $60^{\circ}$}\\
\vect{KL}\cdot\vect{KM} &= \res{32} &&
\end{calculs}
\explique{$\vect{KL}$ et $\vect{LM}$ ne partent pas du même point : on écrit $\vect{KL}=-\vect{LK}$.}{la définition demande deux vecteurs de même origine}
\begin{calculs}
\vect{KL}\cdot\vect{LM} &= -\vect{LK}\cdot\vect{LM} &&\\
\vect{KL}\cdot\vect{LM} &= -8\times 8\times\cos(60^{\circ}) &&\justif{angle $\widehat{KLM}$}\\
\vect{KL}\cdot\vect{LM} &= \res{-32} &&
\end{calculs}
\rem{Erreur fréquente : prendre $60^{\circ}$ entre $\vect{KL}$ et $\vect{LM}$ sans les ramener à la même origine. Mis bout à bout, ils font un angle de $120^{\circ}$.}
\end{exo}

\begin{exo}[Vrai ou faux]
\cas{a)}
\explique{\res{Vrai} : les normes sont positives, donc le cosinus de l'angle est négatif, et l'angle est obtus.}{$\vec{u}\cdot\vec{v}=\norme{\vec{u}}\times\norme{\vec{v}}\times\cos(\alpha)$}
\cas{b)}
\explique{\res{Faux} : dans un carré $ABCD$, $\vect{AB}\cdot\vect{AD}=0$, alors qu'aucun des deux vecteurs n'est nul.}{un produit scalaire nul dit que les vecteurs sont orthogonaux}
\cas{c)}
\explique{\res{Vrai} : l'angle de $\vec{u}$ avec lui-même est nul et $\cos(0)=1$, donc $\vec{u}\cdot\vec{u}=\norme{\vec{u}}\times\norme{\vec{u}}$.}{c'est le carré scalaire}
\cas{d)}
\explique{\res{Faux} : le produit dépend de l'angle ; par exemple $\vec{u}\cdot(-\vec{u})=-9$ quand $\norme{\vec{u}}=3$.}{un contre-exemple suffit}
\end{exo}

\begin{exo}[Hexagone régulier]
\explique{Tous les segments $[OA]$, $[OB]$\dots{} mesurent $2$, et chaque angle au centre $\widehat{AOB}$, $\widehat{BOC}$\dots{} mesure $60^{\circ}$.}{triangles équilatéraux}
\begin{calculs}
\vect{OA}\cdot\vect{OB} &= 2\times 2\times\cos(60^{\circ}) &&\\
\vect{OA}\cdot\vect{OB} &= \res{2} &&\\
\vect{OA}\cdot\vect{OC} &= 2\times 2\times\cos(120^{\circ}) &&\justif{deux angles de $60^{\circ}$}\\
\vect{OA}\cdot\vect{OC} &= \res{-2} &&\justif{$4\times\left(-\frac{1}{2}\right)$}\\
\vect{OA}\cdot\vect{OD} &= 2\times 2\times\cos(180^{\circ}) &&\justif{$A$, $O$ et $D$ sont alignés}\\
\vect{OA}\cdot\vect{OD} &= \res{-4} &&\justif{$\cos(180^{\circ})=-1$}
\end{calculs}
\explique{En $A$, l'angle $\widehat{BAF}$ de l'hexagone vaut $120^{\circ}$ : deux angles de triangles équilatéraux, $\widehat{BAO}$ et $\widehat{OAF}$.}{}
\begin{calculs}
\vect{AB}\cdot\vect{AF} &= 2\times 2\times\cos(120^{\circ}) &&\\
\vect{AB}\cdot\vect{AF} &= \res{-2} &&
\end{calculs}
\end{exo}

\begin{exo}[Travail d'une force]
\cas{a)}
\begin{calculs}
W &= F\times AB\times\cos(\alpha) &&\justif{définition du produit scalaire}\\
W &= 40\times 15\times\cos(30^{\circ}) &&\\
W &= 600\times\dfrac{\sqrt{3}}{2} &&\\
W &= \res{300\sqrt{3}}\text{ J} &&
\end{calculs}
\explique{Soit environ $\res{520}$ J.}{$300\sqrt{3}\approx 519{,}6$}
\cas{b)}
\begin{calculs}
W &= 200\times 300\times\cos(20^{\circ}) &&\\
W &\approx \res{56\,382}\text{ J} &&\justif{calculatrice en mode degré}
\end{calculs}
\cas{c)}
\explique{Le poids est perpendiculaire au déplacement : $\cos(90^{\circ})=0$, donc $\res{W=0}$ J.}{une force perpendiculaire au déplacement ne travaille pas}
\cas{d)}
\begin{calculs}
W &= 30\times 15\times\cos(180^{\circ}) &&\justif{sens contraires : angle de $180^{\circ}$}\\
W &= \res{-450}\text{ J} &&
\end{calculs}
\explique{Le travail est négatif : le frottement \res{freine} le palet.}{travail « résistant »}
\end{exo}

\partie{Produit scalaire et projeté orthogonal}

\begin{exo}[Dans un rectangle]
\cas{a)}
\explique{Le projeté orthogonal de $C$ sur $(AB)$ est $B$.}{$(BC)$ est perpendiculaire à $(AB)$}
\begin{calculs}
\vect{AB}\cdot\vect{AC} &= \vect{AB}\cdot\vect{AB} &&\justif{théorème 3}\\
\vect{AB}\cdot\vect{AC} &= 6\times 6 &&\justif{même sens}\\
\vect{AB}\cdot\vect{AC} &= \res{36} &&
\end{calculs}
\cas{b)}
\explique{$(AB)$ et $(AD)$ sont perpendiculaires : $\res{\vect{AB}\cdot\vect{AD}=0}$.}{vecteurs orthogonaux}
\cas{c)}
\explique{$\vect{CD}=-\vect{AB}$, donc $\vect{AB}\cdot\vect{CD}=-AB\times AB$, soit $\res{-36}$.}{même direction, sens contraires}
\cas{d)}
\explique{Le projeté orthogonal de $I$ sur $(AB)$ est le milieu $I'$ de $[AB]$, et $AI'=3$.}{$I$ est au milieu de $[CD]$, au-dessus du milieu de $[AB]$}
\begin{calculs}
\vect{AB}\cdot\vect{AI} &= AB\times AI' &&\justif{même sens}\\
\vect{AB}\cdot\vect{AI} &= 6\times 3 &&\\
\vect{AB}\cdot\vect{AI} &= \res{18} &&
\end{calculs}
\cas{e)}
\explique{Le projeté orthogonal de $J$ sur $(AD)$ est le milieu $J'$ de $[AD]$, et $AJ'=2$.}{$J$ est à mi-hauteur}
\begin{calculs}
\vect{AD}\cdot\vect{AJ} &= AD\times AJ' &&\justif{même sens}\\
\vect{AD}\cdot\vect{AJ} &= 4\times 2 &&\\
\vect{AD}\cdot\vect{AJ} &= \res{8} &&
\end{calculs}
\cas{f)}
\explique{Le projeté orthogonal de $G$ sur $(EF)$ est $H$ ; comme $E$ est entre $H$ et $F$, les vecteurs $\vect{EF}$ et $\vect{EH}$ sont de sens contraires.}{on fait un petit schéma : $H$, puis $E$, puis $F$ sur la droite}
\begin{calculs}
\vect{EF}\cdot\vect{EG} &= -EF\times EH &&\justif{théorème 3, sens contraires}\\
\vect{EF}\cdot\vect{EG} &= -7\times 2 &&\\
\vect{EF}\cdot\vect{EG} &= \res{-14} &&
\end{calculs}
\end{exo}

\begin{exo}[Sur une droite graduée]
\explique{Les points sont alignés : chaque produit vaut le produit des longueurs, avec le signe « $+$ » si les vecteurs ont le même sens, « $-$ » sinon.}{chaque point est son propre projeté}
\begin{calculs}
\vect{AB}\cdot\vect{AD} &= 2\times 5 &&\justif{$AB=2$, $AD=5$, même sens}\\
\vect{AB}\cdot\vect{AD} &= \res{10} &&\\
\vect{BC}\cdot\vect{BE} &= 1\times 4 &&\justif{même sens}\\
\vect{BC}\cdot\vect{BE} &= \res{4} &&\\
\vect{CA}\cdot\vect{CD} &= -3\times 2 &&\justif{$A$ à gauche de $C$, $D$ à droite : sens contraires}\\
\vect{CA}\cdot\vect{CD} &= \res{-6} &&\\
\vect{AE}\cdot\vect{DB} &= -6\times 3 &&\justif{$\vect{AE}$ va vers la droite, $\vect{DB}$ vers la gauche}\\
\vect{AE}\cdot\vect{DB} &= \res{-18} &&
\end{calculs}
\rem{Le théorème 3 demande deux vecteurs de même origine ; pour deux vecteurs portés par une même droite, la règle « produit des longueurs, avec le signe du sens » vaut toujours.}
\end{exo}

\begin{exo}[Dans un carré]
\cas{1)}
\explique{Le projeté orthogonal de $J$ sur $(AB)$ est $B$.}{$J$ est sur $[BC]$, perpendiculaire à $(AB)$}
\begin{calculs}
\vect{AB}\cdot\vect{AJ} &= AB\times AB &&\justif{même sens}\\
\vect{AB}\cdot\vect{AJ} &= \res{16} &&\\
\vect{AI}\cdot\vect{AJ} &= AI\times AB &&\justif{on projette encore $J$ sur $(AI)$, la même droite}\\
\vect{AI}\cdot\vect{AJ} &= 2\times 4 &&\\
\vect{AI}\cdot\vect{AJ} &= \res{8} &&
\end{calculs}
\cas{2)}
\explique{Le projeté orthogonal de $J$ sur $(DA)$ est le milieu $J'$ de $[DA]$, et $DJ'=2$.}{$J$ est à mi-hauteur du carré}
\begin{calculs}
\vect{DA}\cdot\vect{DJ} &= DA\times DJ' &&\justif{même sens}\\
\vect{DA}\cdot\vect{DJ} &= 4\times 2 &&\\
\vect{DA}\cdot\vect{DJ} &= \res{8} &&
\end{calculs}
\cas{3)}
\begin{calculs}
\vect{IJ} &= \vect{IB}+\vect{BJ} &&\justif{relation de Chasles}\\
\vect{IJ} &= \frac{1}{2}\vect{AB}+\frac{1}{2}\vect{BC} &&\justif{$I$ et $J$ sont des milieux}\\
\vect{IJ} &= \frac{1}{2}\left(\vect{AB}+\vect{BC}\right) &&\justif{on factorise}\\
\vect{IJ} &= \res{\frac{1}{2}\vect{AC}} &&\justif{relation de Chasles}
\end{calculs}
\end{exo}

\partie{Règles de calcul}

\begin{exo}[Calculer avec les règles]
\cas{a)}
\begin{calculs}
(4\vec{u})\cdot(-\vec{v}) &= 4\times(-1)\times\vec{u}\cdot\vec{v} &&\justif{on sort les nombres}\\
(4\vec{u})\cdot(-\vec{v}) &= -4\times 3 &&\\
(4\vec{u})\cdot(-\vec{v}) &= \res{-12} &&
\end{calculs}
\cas{b)}
\begin{calculs}
\vec{u}\cdot(\vec{u}-\vec{v}) &= \vec{u}\cdot\vec{u}-\vec{u}\cdot\vec{v} &&\justif{on développe}\\
\vec{u}\cdot(\vec{u}-\vec{v}) &= \norme{\vec{u}}^{2}-\vec{u}\cdot\vec{v} &&\justif{carré scalaire}\\
\vec{u}\cdot(\vec{u}-\vec{v}) &= 4-3 &&\\
\vec{u}\cdot(\vec{u}-\vec{v}) &= \res{1} &&
\end{calculs}
\cas{c)}
\begin{calculs}
\norme{\vec{u}+\vec{v}}^{2} &= \norme{\vec{u}}^{2}+2\,\vec{u}\cdot\vec{v}+\norme{\vec{v}}^{2} &&\justif{première identité remarquable}\\
\norme{\vec{u}+\vec{v}}^{2} &= 4+6+25 &&\\
\norme{\vec{u}+\vec{v}}^{2} &= 35 &&\\
\norme{\vec{u}+\vec{v}} &= \res{\sqrt{35}} &&\justif{une norme est positive}
\end{calculs}
\cas{d)}
\begin{calculs}
\norme{\vec{u}-\vec{v}}^{2} &= \norme{\vec{u}}^{2}-2\,\vec{u}\cdot\vec{v}+\norme{\vec{v}}^{2} &&\justif{deuxième identité remarquable}\\
\norme{\vec{u}-\vec{v}}^{2} &= 4-6+25 &&\\
\norme{\vec{u}-\vec{v}}^{2} &= 23 &&\\
\norme{\vec{u}-\vec{v}} &= \res{\sqrt{23}} &&
\end{calculs}
\cas{e)}
\begin{calculs}
(\vec{u}+2\vec{v})\cdot(\vec{u}-\vec{v}) &= \vec{u}\cdot\vec{u}-\vec{u}\cdot\vec{v}+2\vec{v}\cdot\vec{u}-2\vec{v}\cdot\vec{v} &&\justif{on développe, comme $(a+2b)(a-b)$}\\
(\vec{u}+2\vec{v})\cdot(\vec{u}-\vec{v}) &= \norme{\vec{u}}^{2}+\vec{u}\cdot\vec{v}-2\norme{\vec{v}}^{2} &&\justif{$-\vec{u}\cdot\vec{v}+2\vec{u}\cdot\vec{v}$ : symétrie}\\
(\vec{u}+2\vec{v})\cdot(\vec{u}-\vec{v}) &= 4+3-50 &&\\
(\vec{u}+2\vec{v})\cdot(\vec{u}-\vec{v}) &= \res{-43} &&
\end{calculs}
\end{exo}

\begin{exo}[Diagonales d'un losange]
\cas{1)}
\begin{calculs}
(\vec{u}+\vec{v})\cdot(\vec{u}-\vec{v}) &= \vec{u}\cdot\vec{u}-\vec{u}\cdot\vec{v}+\vec{v}\cdot\vec{u}-\vec{v}\cdot\vec{v} &&\justif{on développe}\\
(\vec{u}+\vec{v})\cdot(\vec{u}-\vec{v}) &= \vec{u}\cdot\vec{u}-\vec{v}\cdot\vec{v} &&\justif{$\vec{v}\cdot\vec{u}=\vec{u}\cdot\vec{v}$ : les termes s'annulent}\\
(\vec{u}+\vec{v})\cdot(\vec{u}-\vec{v}) &= \res{\norme{\vec{u}}^{2}-\norme{\vec{v}}^{2}} &&\justif{carrés scalaires}
\end{calculs}
\cas{2)}
\begin{calculs}
\vect{AC} &= \vect{AB}+\vect{BC} &&\justif{relation de Chasles}\\
\vect{AC} &= \res{\vect{AB}+\vect{AD}} &&\justif{parallélogramme : $\vect{BC}=\vect{AD}$}\\
\vect{DB} &= \vect{DA}+\vect{AB} &&\justif{relation de Chasles}\\
\vect{DB} &= \res{\vect{AB}-\vect{AD}} &&\justif{$\vect{DA}=-\vect{AD}$}
\end{calculs}
\cas{3)}
\begin{calculs}
\vect{AC}\cdot\vect{DB} &= (\vect{AB}+\vect{AD})\cdot(\vect{AB}-\vect{AD}) &&\justif{question 2}\\
\vect{AC}\cdot\vect{DB} &= AB^{2}-AD^{2} &&\justif{question 1}\\
\vect{AC}\cdot\vect{DB} &= 0 &&\justif{losange : $AB=AD$}
\end{calculs}
\explique{Le produit scalaire est nul : les diagonales d'un losange \res{sont perpendiculaires}.}{}
\end{exo}

\begin{exo}[Angle obtus dans un rectangle]
\cas{1)}
\explique{$\res{\vect{KD}=\vect{KA}+\vect{AD}}$ et $\res{\vect{KL}=\vect{KB}+\vect{BL}}$.}{relation de Chasles}
\cas{2)}
\explique{$\vect{KA}$ est porté par $(AB)$ et $\vect{BL}$ par $(BC)$, perpendiculaire à $(AB)$ : $\res{\vect{KA}\cdot\vect{BL}=0}$.}{angle droit en $B$}
\explique{$\vect{AD}$ est porté par $(AD)$ et $\vect{KB}$ par $(AB)$ : $\res{\vect{AD}\cdot\vect{KB}=0}$.}{angle droit en $A$}
\cas{3)}
\explique{$K$ est entre $A$ et $B$, avec $KA=4$ et $KB=2$ : $\vect{KA}$ et $\vect{KB}$ sont de sens contraires, donc $\vect{KA}\cdot\vect{KB}=-4\times 2$, soit $\res{-8}$.}{$KB=6-4$}
\explique{$\vect{AD}$ et $\vect{BL}$ ont la même direction et le même sens (vers le haut) : $\vect{AD}\cdot\vect{BL}=4\times 1$, soit $\res{4}$.}{}
\cas{4)}
\begin{calculs}
\vect{KD}\cdot\vect{KL} &= (\vect{KA}+\vect{AD})\cdot(\vect{KB}+\vect{BL}) &&\justif{question 1}\\
\vect{KD}\cdot\vect{KL} &= \vect{KA}\cdot\vect{KB}+\vect{KA}\cdot\vect{BL}+\vect{AD}\cdot\vect{KB}+\vect{AD}\cdot\vect{BL} &&\justif{on développe}\\
\vect{KD}\cdot\vect{KL} &= -8+0+0+4 &&\justif{questions 2 et 3}\\
\vect{KD}\cdot\vect{KL} &= \res{-4} &&
\end{calculs}
\explique{Le produit scalaire est négatif : l'angle $\widehat{DKL}$ \res{est obtus}.}{propriété 7}
\rem{Vérification dans le repère d'origine $A$ : $K(4\,;0)$, $D(0\,;4)$, $L(6\,;1)$, $\vect{KD}\coord{-4}{4}$ et $\vect{KL}\coord{2}{1}$ ; on trouve $-8+4$, soit $-4$.}
\end{exo}

\partie{Produit scalaire dans un repère orthonormé}

\begin{exo}[Produit scalaire et orthogonalité]
\cas{a)}
\begin{calculs}
\vec{u}\cdot\vec{v} &= 5\times 3+(-2)\times 4 &&\justif{$xx'+yy'$}\\
\vec{u}\cdot\vec{v} &= \res{7} &&\justif{$15-8$}\\
\vec{u}\,{}'\cdot\vec{v}\,{}' &= \dfrac{1}{3}\times 6+2\times\left(-\dfrac{1}{2}\right) &&\\
\vec{u}\,{}'\cdot\vec{v}\,{}' &= 2-1 &&\justif{$\frac{6}{3}=2$}\\
\vec{u}\,{}'\cdot\vec{v}\,{}' &= \res{1} &&
\end{calculs}
\cas{b)}
\begin{calculs}
\vec{a}\cdot\vec{b} &= -4\times 3+6\times 2 &&\\
\vec{a}\cdot\vec{b} &= 0 &&\justif{$-12+12$}
\end{calculs}
\explique{Le produit scalaire est nul : $\vec{a}$ et $\vec{b}$ \res{sont orthogonaux}.}{}
\cas{c)}
\begin{calculs}
\vect{AB} & \coord{2}{4} &&\justif{$4-2$ et $3-(-1)$}\\
\vect{AC} & \coord{-4}{2} &&\justif{$-2-2$ et $1-(-1)$}\\
\vect{AB}\cdot\vect{AC} &= 2\times(-4)+4\times 2 &&\\
\vect{AB}\cdot\vect{AC} &= 0 &&
\end{calculs}
\explique{$\vect{AB}$ et $\vect{AC}$ sont orthogonaux : le triangle $ABC$ \res{est rectangle en $A$}.}{deux vecteurs issus de $A$}
\cas{d)}
\begin{calculs}
\vec{r}\cdot\vec{s} &= 4m+2\times(-6) &&\\
\vec{r}\cdot\vec{s} &= 4m-12 &&\\
4m-12 &= 0 &&\justif{orthogonaux : produit scalaire nul}\\
m &= 3 &&
\end{calculs}
\explique{$\vec{r}$ et $\vec{s}$ sont orthogonaux pour $\res{m=3}$.}{}
\end{exo}

\begin{exo}[Rectangle en quel sommet ?]
\cas{a)}
\begin{calculs}
\vect{KL} & \coord{4}{2} &&\justif{on essaie d'abord le sommet $K$}\\
\vect{KM} & \coord{-1}{2} &&\\
\vect{KL}\cdot\vect{KM} &= 4\times(-1)+2\times 2 &&\\
\vect{KL}\cdot\vect{KM} &= 0 &&
\end{calculs}
\explique{Le triangle $KLM$ \res{est rectangle en $K$}.}{un triangle a au plus un angle droit : inutile de tester $L$ et $M$}
\cas{b)}
\begin{calculs}
\vect{RS} & \coord{4}{2} &&\justif{sommet $R$}\\
\vect{RT} & \coord{6}{-2} &&\\
\vect{RS}\cdot\vect{RT} &= 24-4 &&\\
\vect{RS}\cdot\vect{RT} &= 20 &&\justif{pas d'angle droit en $R$}\\
\vect{SR} & \coord{-4}{-2} &&\justif{sommet $S$}\\
\vect{ST} & \coord{2}{-4} &&\\
\vect{SR}\cdot\vect{ST} &= -8+8 &&\\
\vect{SR}\cdot\vect{ST} &= 0 &&
\end{calculs}
\explique{Le triangle $RST$ \res{est rectangle en $S$}.}{}
\end{exo}

\begin{exo}[Vecteurs de la base]
\cas{1)}
\begin{calculs}
\vec{u}\cdot\vec{\imath} &= -2\times 1+7\times 0 &&\justif{$\vec{\imath}\coord{1}{0}$}\\
\vec{u}\cdot\vec{\imath} &= \res{-2} &&\justif{l'abscisse de $\vec{u}$}\\
\vec{u}\cdot\vec{\jmath} &= -2\times 0+7\times 1 &&\justif{$\vec{\jmath}\coord{0}{1}$}\\
\vec{u}\cdot\vec{\jmath} &= \res{7} &&\justif{l'ordonnée de $\vec{u}$}
\end{calculs}
\cas{2)}
\explique{$\vec{w}\cdot\vec{\imath}$ est l'abscisse de $\vec{w}$ et $\vec{w}\cdot\vec{\jmath}$ son ordonnée : $\res{\vec{w}\coord{4}{-1}}$.}{question 1}
\cas{3)}
\begin{calculs}
\vec{u}\cdot\vec{v} &= -2\times 7+7\times 2 &&\\
\vec{u}\cdot\vec{v} &= 0 &&
\end{calculs}
\explique{$\vec{v}$ \res{est orthogonal} à $\vec{u}$.}{}
\cas{4)}
\begin{calculs}
\vec{u}\cdot\vec{t} &= -2\times 14+7\times 4 &&\\
\vec{u}\cdot\vec{t} &= 0 &&\justif{$-28+28$}
\end{calculs}
\explique{$\vec{t}$ \res{est orthogonal} à $\vec{u}$ ; comme $\vec{t}=2\vec{v}$, il \res{est colinéaire} à $\vec{v}$.}{$\det(\vec{t}\,;\vec{v})=14\times 2-4\times 7=0$}
\end{exo}

\begin{exo}[Où placer le point ?]
\cas{1)}
\begin{calculs}
\vect{AB} & \coord{4}{-2} &&\justif{$5-1$ et $0-2$}\\
AB &= \sqrt{16+4} &&\\
AB &= \sqrt{20} &&\\
AB &= \res{2\sqrt{5}} &&\justif{$\sqrt{4}\times\sqrt{5}$}
\end{calculs}
\cas{2)}
\begin{calculs}
\vect{AM} & \coord{x-1}{1} &&\justif{$3-2$}\\
\vect{AB}\cdot\vect{AM} &= 4(x-1)+(-2)\times 1 &&\\
\vect{AB}\cdot\vect{AM} &= 4x-6 &&\\
4x-6 &= 0 &&\justif{rectangle en $A$ : produit scalaire nul}\\
x &= \dfrac{3}{2} &&
\end{calculs}
\explique{Le triangle $ABM$ est rectangle en $A$ pour $\res{x=\frac{3}{2}}$.}{}
\cas{3)}
\begin{calculs}
\vect{MA} & \coord{1-x}{-1} &&\justif{arrivée moins départ : $2-3$}\\
\vect{MB} & \coord{5-x}{-3} &&\justif{$0-3$}\\
\vect{MA}\cdot\vect{MB} &= (1-x)(5-x)+(-1)\times(-3) &&\\
\vect{MA}\cdot\vect{MB} &= 5-x-5x+x^{2}+3 &&\justif{on développe}\\
\vect{MA}\cdot\vect{MB} &= x^{2}-6x+8 &&
\end{calculs}
\explique{Le triangle $ABM$ est rectangle en $M$ quand $\vect{MA}\cdot\vect{MB}=0$, c'est-à-dire \res{quand $x^{2}-6x+8=0$}.}{}
\cas{4)}
\begin{calculs}
\Delta &= (-6)^{2}-4\times 1\times 8 &&\justif{discriminant, chapitre 3}\\
\Delta &= 4 &&\\
x_{1} &= \dfrac{6-2}{2} &&\\
x_{1} &= 2 &&\\
x_{2} &= \dfrac{6+2}{2} &&\\
x_{2} &= 4 &&
\end{calculs}
\explique{Les points cherchés sont $\res{M_{1}(2\,;3)}$ et $\res{M_{2}(4\,;3)}$.}{on n'oublie pas l'ordonnée $3$}
\end{exo}

\partie{Produit scalaire et normes}

\begin{exo}[Produit scalaire avec des normes]
\cas{a)}
\begin{calculs}
\vec{u}\cdot\vec{v} &= \frac{1}{2}\left(\norme{\vec{u}+\vec{v}}^{2}-\norme{\vec{u}}^{2}-\norme{\vec{v}}^{2}\right) &&\justif{on recopie la formule}\\
\vec{u}\cdot\vec{v} &= \frac{1}{2}\left(64-16-36\right) &&\justif{$8^{2}$, $4^{2}$ et $6^{2}$}\\
\vec{u}\cdot\vec{v} &= \res{6} &&
\end{calculs}
\cas{b)}
\begin{calculs}
\vect{AB}\cdot\vect{AC} &= \frac{1}{2}\left(AB^{2}+AC^{2}-BC^{2}\right) &&\justif{sommet commun $A$ : on retranche $BC^{2}$}\\
\vect{AB}\cdot\vect{AC} &= \frac{1}{2}\left(25+64-49\right) &&\\
\vect{AB}\cdot\vect{AC} &= \res{20} &&\\
\vect{BA}\cdot\vect{BC} &= \frac{1}{2}\left(BA^{2}+BC^{2}-AC^{2}\right) &&\justif{sommet commun $B$}\\
\vect{BA}\cdot\vect{BC} &= \frac{1}{2}\left(25+49-64\right) &&\\
\vect{BA}\cdot\vect{BC} &= \res{5} &&\\
\vect{CA}\cdot\vect{CB} &= \frac{1}{2}\left(CA^{2}+CB^{2}-AB^{2}\right) &&\justif{sommet commun $C$}\\
\vect{CA}\cdot\vect{CB} &= \frac{1}{2}\left(64+49-25\right) &&\\
\vect{CA}\cdot\vect{CB} &= \res{44} &&
\end{calculs}
\cas{c)}
\begin{calculs}
\norme{\vec{u}-\vec{v}}^{2} &= \norme{\vec{u}}^{2}-2\,\vec{u}\cdot\vec{v}+\norme{\vec{v}}^{2} &&\justif{deuxième identité remarquable}\\
16 &= 9-2\,\vec{u}\cdot\vec{v}+4 &&\justif{$4^{2}$, $3^{2}$ et $2^{2}$}\\
2\,\vec{u}\cdot\vec{v} &= -3 &&\justif{$13-16$}\\
\vec{u}\cdot\vec{v} &= \res{-\frac{3}{2}} &&
\end{calculs}
\end{exo}

\begin{exo}[Diagonales d'un parallélogramme]
\cas{1)}
\begin{calculs}
\vect{AB}\cdot\vect{AD} &= \frac{1}{2}\left(AB^{2}+AD^{2}-BD^{2}\right) &&\justif{triangle $ABD$, sommet commun $A$}\\
\vect{AB}\cdot\vect{AD} &= \frac{1}{2}\left(36+16-25\right) &&\\
\vect{AB}\cdot\vect{AD} &= \res{\frac{27}{2}} &&
\end{calculs}
\cas{2)}
\begin{calculs}
AC^{2} &= \norme{\vect{AB}+\vect{AD}}^{2} &&\justif{règle du parallélogramme}\\
AC^{2} &= AB^{2}+2\,\vect{AB}\cdot\vect{AD}+AD^{2} &&\justif{première identité remarquable}\\
AC^{2} &= 36+27+16 &&\justif{$2\times\frac{27}{2}=27$}\\
AC^{2} &= \res{79} &&\\
AC &= \res{\sqrt{79}} &&
\end{calculs}
\cas{3)}
\begin{calculs}
AC^{2}+BD^{2} &= 79+25 &&\\
AC^{2}+BD^{2} &= 104 &&\\
2\left(AB^{2}+AD^{2}\right) &= 2\times(36+16) &&\\
2\left(AB^{2}+AD^{2}\right) &= 104 &&
\end{calculs}
\explique{Les deux membres valent $104$ : \res{l'égalité est vérifiée}.}{la somme des carrés des diagonales vaut la somme des carrés des quatre côtés}
\end{exo}

\begin{exo}[Une voile triangulaire]
\cas{1)}
\begin{calculs}
\vect{AB}\cdot\vect{AC} &= \frac{1}{2}\left(AB^{2}+AC^{2}-BC^{2}\right) &&\justif{sommet commun $A$}\\
\vect{AB}\cdot\vect{AC} &= \frac{1}{2}\left(9+16-36\right) &&\\
\vect{AB}\cdot\vect{AC} &= \res{-\frac{11}{2}} &&\\
\vect{BA}\cdot\vect{BC} &= \frac{1}{2}\left(9+36-16\right) &&\justif{sommet commun $B$ : on retranche $AC^{2}$}\\
\vect{BA}\cdot\vect{BC} &= \res{\frac{29}{2}} &&\\
\vect{CA}\cdot\vect{CB} &= \frac{1}{2}\left(16+36-9\right) &&\justif{sommet commun $C$ : on retranche $AB^{2}$}\\
\vect{CA}\cdot\vect{CB} &= \res{\frac{43}{2}} &&
\end{calculs}
\cas{2)}
\explique{L'angle en $A$ \res{est obtus}, car $\vect{AB}\cdot\vect{AC}$ est négatif ; les angles en $B$ et en $C$ \res{sont aigus}, car les deux autres produits sont positifs.}{signe du produit scalaire, propriété 7}
\end{exo}

\partie{Vecteur normal et équation cartésienne d'une droite}

\begin{exo}[Équation avec un vecteur normal]
\cas{a)}
\explique{$(d_{1})$ : $5x+(-3)y+2=0$ ; $\res{\vec{n}_{1}\coord{5}{-3}}$ est normal, $\res{\vec{u}_{1}\coord{3}{5}}$ est directeur.}{$\vec{n}\coord{a}{b}$ et $\vec{u}\coord{-b}{a}$}
\explique{$(d_{2})$ : $-2x+1y+0=0$ ; $\res{\vec{n}_{2}\coord{-2}{1}}$ est normal, $\res{\vec{u}_{2}\coord{-1}{-2}}$ est directeur.}{pas de $c$ : la droite passe par l'origine}
\explique{$(d_{3})$ s'écrit $4x-y-1=0$ ; $\res{\vec{n}_{3}\coord{4}{-1}}$ est normal, $\res{\vec{u}_{3}\coord{1}{4}}$ est directeur.}{on passe $y$ dans l'autre membre ; on retrouve le coefficient directeur $4$}
\cas{b)}
\begin{calculs}
2x+5y+c &= 0 &&\justif{$a=2$ et $b=5$, lus sur $\vec{n}$}\\
2\times 1+5\times(-3)+c &= 0 &&\justif{$A(1\,;-3)$ est sur la droite}\\
-13+c &= 0 &&\\
c &= 13 &&
\end{calculs}
\explique{Une équation de la droite est $\res{2x+5y+13=0}$.}{on remplace $c$}
\cas{c)}
\explique{La hauteur issue de $B$ passe par $B$ et elle est perpendiculaire à $(CD)$ : $\vect{CD}$ est un vecteur normal.}{on cherche d'abord le vecteur normal}
\begin{calculs}
\vect{CD} & \coord{4}{-5} &&\justif{$3-(-1)$ et $-2-3$}\\
4x-5y+c &= 0 &&\justif{$a=4$ et $b=-5$}\\
4\times 2-5\times 1+c &= 0 &&\justif{$B(2\,;1)$ est sur la hauteur}\\
c &= -3 &&\justif{$3+c=0$}
\end{calculs}
\explique{Une équation de la hauteur est $\res{4x-5y-3=0}$.}{}
\cas{d)}
\explique{La médiatrice de $[EF]$ passe par le milieu $I$ de $[EF]$ et elle est perpendiculaire à $(EF)$ : $\vect{EF}$ est un vecteur normal.}{définition de la médiatrice}
\begin{calculs}
I & \left(\dfrac{-3+1}{2}\,;\dfrac{2+6}{2}\right) &&\justif{coordonnées du milieu}\\
I & (-1\,;4) &&\\
\vect{EF} & \coord{4}{4} &&\justif{$1-(-3)$ et $6-2$}\\
4x+4y+c &= 0 &&\\
4\times(-1)+4\times 4+c &= 0 &&\justif{$I$ est sur la médiatrice}\\
c &= -12 &&\justif{$12+c=0$}
\end{calculs}
\explique{Une équation de la médiatrice est $\res{4x+4y-12=0}$, ou encore $x+y-3=0$.}{on peut diviser les deux membres par $4$}
\end{exo}

\begin{exo}[Parallèles ou perpendiculaires ?]
\cas{a)}
\begin{calculs}
\vec{n}_{1} & \coord{3}{4} &&\justif{normal à la première droite}\\
\vec{n}_{2} & \coord{4}{-3} &&\justif{normal à la seconde}\\
\vec{n}_{1}\cdot\vec{n}_{2} &= 3\times 4+4\times(-3) &&\justif{$xx'+yy'$}\\
\vec{n}_{1}\cdot\vec{n}_{2} &= 0 &&
\end{calculs}
\explique{Les vecteurs normaux sont orthogonaux : les droites \res{sont perpendiculaires}.}{propriété 11}
\cas{b)}
\begin{calculs}
\vec{n}_{3} & \coord{-6}{-8} &&\\
\det(\vec{n}_{1}\,;\vec{n}_{3}) &= 3\times(-8)-4\times(-6) &&\justif{$xy'-yx'$}\\
\det(\vec{n}_{1}\,;\vec{n}_{3}) &= 0 &&\justif{$-24+24$}
\end{calculs}
\explique{Les vecteurs normaux sont colinéaires : les droites \res{sont parallèles}.}{d'ailleurs $\vec{n}_{3}=-2\,\vec{n}_{1}$}
\cas{c)}
\explique{On écrit la seconde droite sous forme cartésienne : $3x-y+1=0$.}{on passe $y$ dans l'autre membre}
\begin{calculs}
\vec{n}_{4} & \coord{1}{-2} &&\\
\vec{n}_{5} & \coord{3}{-1} &&\\
\vec{n}_{4}\cdot\vec{n}_{5} &= 3+2 &&\\
\vec{n}_{4}\cdot\vec{n}_{5} &= 5 &&\justif{non nul : pas perpendiculaires}\\
\det(\vec{n}_{4}\,;\vec{n}_{5}) &= 1\times(-1)-(-2)\times 3 &&\\
\det(\vec{n}_{4}\,;\vec{n}_{5}) &= 5 &&\justif{non nul : pas parallèles}
\end{calculs}
\explique{Les droites ne sont \res{ni parallèles ni perpendiculaires}.}{}
\cas{d)}
\explique{$x=4$ s'écrit $1x+0y-4=0$ et $y=-2$ s'écrit $0x+1y+2=0$ : $\vec{n}\coord{1}{0}$ et $\vec{n}\,{}'\coord{0}{1}$ sont orthogonaux ; les droites \res{sont perpendiculaires}.}{une droite verticale et une droite horizontale}
\end{exo}

\begin{exo}[Tangentes à un cercle]
\cas{1)}
\explique{$\mathcal{C}$ : $\res{(x-1)^{2}+(y+2)^{2}=25}$.}{centre $\Omega(1\,;-2)$, rayon $5$}
\begin{calculs}
(4-1)^{2}+(2+2)^{2} &= 9+16 &&\justif{coordonnées de $A$}\\
(4-1)^{2}+(2+2)^{2} &= 25 &&\\
(-2-1)^{2}+(2+2)^{2} &= 9+16 &&\justif{coordonnées de $B$}\\
(-2-1)^{2}+(2+2)^{2} &= 25 &&
\end{calculs}
\explique{$A$ et $B$ \res{sont sur $\mathcal{C}$}.}{}
\cas{2)}
\explique{La tangente en $A$ est perpendiculaire à $(\Omega A)$ : $\vect{\Omega A}$ est un vecteur normal.}{}
\begin{calculs}
\vect{\Omega A} & \coord{3}{4} &&\justif{$4-1$ et $2-(-2)$}\\
3x+4y+c &= 0 &&\\
3\times 4+4\times 2+c &= 0 &&\justif{$A(4\,;2)$ est sur la tangente}\\
c &= -20 &&
\end{calculs}
\explique{Tangente en $A$ : $\res{3x+4y-20=0}$.}{}
\begin{calculs}
\vect{\Omega B} & \coord{-3}{4} &&\justif{$-2-1$ et $2-(-2)$}\\
-3x+4y+c &= 0 &&\\
-3\times(-2)+4\times 2+c &= 0 &&\justif{$B(-2\,;2)$ est sur la tangente}\\
c &= -14 &&
\end{calculs}
\explique{Tangente en $B$ : $\res{-3x+4y-14=0}$.}{}
\cas{3)}
\begin{calculs}
\vect{\Omega A}\cdot\vect{\Omega B} &= 3\times(-3)+4\times 4 &&\justif{produit des vecteurs normaux}\\
\vect{\Omega A}\cdot\vect{\Omega B} &= 7 &&
\end{calculs}
\explique{Le produit n'est pas nul : les deux tangentes \res{ne sont pas perpendiculaires}.}{propriété 11}
\end{exo}

\begin{exo}[Les trois hauteurs]
\cas{1)}
\explique{La hauteur issue de $A$ est perpendiculaire à $(BC)$ : $\vect{BC}$ est normal.}{}
\begin{calculs}
\vect{BC} & \coord{6}{0} &&\justif{$4-(-2)$ et $0-0$}\\
6x+0y+c &= 0 &&\\
6\times 0+c &= 0 &&\justif{$A(0\,;4)$ est sur la hauteur}\\
c &= 0 &&
\end{calculs}
\explique{Hauteur issue de $A$ : $6x=0$, soit $\res{x=0}$.}{l'axe des ordonnées, car $(BC)$ est l'axe des abscisses}
\begin{calculs}
\vect{AC} & \coord{4}{-4} &&\justif{normal à la hauteur issue de $B$}\\
4x-4y+c &= 0 &&\\
4\times(-2)-4\times 0+c &= 0 &&\justif{$B(-2\,;0)$ est sur la hauteur}\\
c &= 8 &&
\end{calculs}
\explique{Hauteur issue de $B$ : $\res{4x-4y+8=0}$, ou $x-y+2=0$.}{on divise par $4$}
\cas{2)}
\explique{$H(x\,;y)$ est solution du système $\systeme{x=0}{x-y+2=0}$.}{on le résout par substitution}
\explique{En remplaçant $x$ par $0$, on obtient $-y+2=0$, donc $y=2$, et $\res{H(0\,;2)}$.}{}
\cas{3)}
\begin{calculs}
\vect{CH} & \coord{-4}{2} &&\justif{$0-4$ et $2-0$}\\
\vect{AB} & \coord{-2}{-4} &&\justif{$-2-0$ et $0-4$}\\
\vect{CH}\cdot\vect{AB} &= 8-8 &&\\
\vect{CH}\cdot\vect{AB} &= 0 &&
\end{calculs}
\explique{$(CH)$ et $(AB)$ \res{sont perpendiculaires} : $(CH)$ est la hauteur issue de $C$. Les trois hauteurs du triangle \res{passent par le même point} $H$.}{elles sont concourantes : $H$ est l'orthocentre}
\cas{4)}
\begin{calculs}
\vect{AB}\cdot\vect{AC} &= -2\times 4+(-4)\times(-4) &&\\
\vect{AB}\cdot\vect{AC} &= \res{8} &&\justif{$-8+16$}
\end{calculs}
\explique{Le produit scalaire est positif : l'angle $\widehat{BAC}$ \res{est aigu}.}{propriété 7}
\end{exo}

\partie{Projeté orthogonal d'un point sur une droite}

\begin{exo}[Projeté orthogonal]
\cas{1) a)}
\explique{$\vec{u}\coord{-2}{1}$ dirige $(d)$, donc il est normal à $(\Delta)$.}{$\vec{u}\coord{-b}{a}$ avec $a=1$ et $b=2$}
\begin{calculs}
-2x+y+c &= 0 &&\justif{$a=-2$ et $b=1$}\\
-2\times 3+6+c &= 0 &&\justif{$M(3\,;6)$ est sur $(\Delta)$}\\
c &= 0 &&
\end{calculs}
\explique{Une équation de $(\Delta)$ est $\res{-2x+y=0}$.}{}
\cas{1) b)}
\explique{$H(x\,;y)$ est solution du système $\systeme{x+2y-5=0}{-2x+y=0}$.}{équations de $(d)$, puis de $(\Delta)$}
\begin{calculs}
y &= 2x &&\justif{on isole $y$ dans la seconde équation}\\
x+2\times 2x-5 &= 0 &&\justif{on remplace $y$ dans la première}\\
5x-5 &= 0 &&\\
x &= 1 &&\\
y &= 2 &&\justif{$y=2\times 1$}
\end{calculs}
\explique{Le projeté orthogonal de $M$ sur $(d)$ est $\res{H(1\,;2)}$.}{on vérifie dans l'équation de $(d)$ : $1+4-5=0$}
\cas{1) c)}
\begin{calculs}
\vect{MH} & \coord{-2}{-4} &&\justif{$1-3$ et $2-6$}\\
MH &= \sqrt{4+16} &&\\
MH &= \sqrt{20} &&\\
MH &= \res{2\sqrt{5}} &&\justif{$\sqrt{4}\times\sqrt{5}$}
\end{calculs}
\explique{La distance du point $M$ à la droite $(d)$ est $2\sqrt{5}$.}{c'est la plus courte distance de $M$ à un point de $(d)$}
\cas{2)}
\explique{On cherche d'abord une équation de $(AB)$ : $\vect{AB}\coord{4}{2}$ la dirige, donc $\vec{n}\coord{1}{-2}$ lui est normal.}{$\vect{AB}\cdot\vec{n}=4-4=0$}
\begin{calculs}
x-2y+c &= 0 &&\justif{équation de $(AB)$}\\
-1-2\times 0+c &= 0 &&\justif{$A(-1\,;0)$ est sur $(AB)$}\\
c &= 1 &&
\end{calculs}
\explique{$(AB)$ : $x-2y+1=0$. La perpendiculaire à $(AB)$ qui passe par $C$ a pour vecteur normal $\vect{AB}\coord{4}{2}$.}{puis la perpendiculaire}
\begin{calculs}
4x+2y+c &= 0 &&\\
4\times 0+2\times 3+c &= 0 &&\justif{$C(0\,;3)$ est sur la perpendiculaire}\\
c &= -6 &&
\end{calculs}
\explique{$K(x\,;y)$ est solution du système $\systeme{x-2y+1=0}{4x+2y-6=0}$ ; on ajoute les deux équations.}{les coefficients de $y$ sont opposés : $-2y$ et $+2y$ s'éliminent}
\begin{calculs}
5x-5 &= 0 &&\justif{$(x+4x)+(-2y+2y)+(1-6)$}\\
x &= 1 &&\\
1-2y+1 &= 0 &&\justif{on remplace $x$ par $1$ dans la première équation}\\
y &= 1 &&\justif{$2-2y=0$}
\end{calculs}
\explique{Le projeté orthogonal de $C$ sur $(AB)$ est $\res{K(1\,;1)}$.}{on vérifie : $\vect{CK}\coord{1}{-2}$ et $\vect{CK}\cdot\vect{AB}=4-4=0$}
\end{exo}

\begin{exo}[Distance d'un point à une droite]
\cas{a)}
\explique{$\vec{u}\coord{-1}{1}$ dirige $(d)$ : la perpendiculaire $(\Delta)$ à $(d)$ qui passe par $A$ a pour équation $-x+y+c=0$.}{$\vec{u}\coord{-b}{a}$}
\begin{calculs}
-5+5+c &= 0 &&\justif{$A(5\,;5)$ est sur $(\Delta)$}\\
c &= 0 &&
\end{calculs}
\explique{$(\Delta)$ : $-x+y=0$, soit $y=x$. Le projeté $H$ est solution de $\systeme{x+y-4=0}{y=x}$.}{substitution}
\begin{calculs}
x+x-4 &= 0 &&\justif{on remplace $y$ par $x$}\\
x &= 2 &&\\
y &= 2 &&
\end{calculs}
\begin{calculs}
\vect{AH} & \coord{-3}{-3} &&\justif{$H(2\,;2)$}\\
AH &= \sqrt{9+9} &&\\
AH &= \res{3\sqrt{2}} &&\justif{$\sqrt{18}=\sqrt{9}\times\sqrt{2}$}
\end{calculs}
\explique{La distance du point $A$ à la droite $(d)$ est $3\sqrt{2}$.}{}
\cas{b)}
\explique{$\vec{u}\coord{4}{3}$ dirige $(d')$ : la perpendiculaire à $(d')$ qui passe par $B$ a pour équation $4x+3y+c=0$.}{$\vec{u}\coord{-b}{a}$ avec $a=3$ et $b=-4$}
\begin{calculs}
4\times(-1)+3\times 9+c &= 0 &&\justif{$B(-1\,;9)$ est sur la perpendiculaire}\\
c &= -23 &&\justif{$23+c=0$}
\end{calculs}
\explique{Le projeté $H$ est solution de $\systeme{3x-4y+14=0}{4x+3y-23=0}$ ; on multiplie la première équation par $3$ et la seconde par $4$.}{combinaison : on fait apparaître $-12y$ et $+12y$}
\begin{calculs}
9x-12y+42 &= 0 &&\justif{première équation fois $3$}\\
16x+12y-92 &= 0 &&\justif{seconde équation fois $4$}\\
25x-50 &= 0 &&\justif{on ajoute les deux équations}\\
x &= 2 &&\\
3\times 2-4y+14 &= 0 &&\justif{on remplace $x$ dans la première équation}\\
y &= 5 &&\justif{$20-4y=0$}
\end{calculs}
\begin{calculs}
\vect{BH} & \coord{3}{-4} &&\justif{$H(2\,;5)$}\\
BH &= \sqrt{9+16} &&\\
BH &= \res{5} &&
\end{calculs}
\explique{La distance du point $B$ à la droite $(d')$ est $5$.}{}
\end{exo}

\begin{exo}[Raccorder une maison]
\cas{1)}
\explique{$\vec{u}\coord{2}{1}$ dirige $(d)$, donc il est normal à la perpendiculaire.}{$\vec{u}\coord{-b}{a}$ avec $a=1$ et $b=-2$}
\begin{calculs}
2x+y+c &= 0 &&\\
2\times 6+2+c &= 0 &&\justif{$M(6\,;2)$ est sur la perpendiculaire}\\
c &= -14 &&
\end{calculs}
\explique{Une équation de la perpendiculaire est $\res{2x+y-14=0}$.}{}
\cas{2)}
\explique{$H$ est solution du système $\systeme{x-2y+8=0}{2x+y-14=0}$ ; on le résout par substitution.}{$x$ s'isole facilement dans la première équation}
\begin{calculs}
x &= 2y-8 &&\justif{on isole $x$}\\
2(2y-8)+y-14 &= 0 &&\justif{on remplace $x$ dans la seconde}\\
5y-30 &= 0 &&\justif{$4y-16+y-14$}\\
y &= 6 &&\\
x &= 4 &&\justif{$x=2\times 6-8$}
\end{calculs}
\explique{Le tuyau arrive au point $\res{H(4\,;6)}$.}{on vérifie dans la seconde équation : $8+6-14=0$}
\cas{3)}
\begin{calculs}
\vect{MH} & \coord{-2}{4} &&\\
MH &= \sqrt{4+16} &&\\
MH &= 2\sqrt{5} &&\justif{en hectomètres}\\
MH &\approx 4{,}472 &&
\end{calculs}
\explique{Le tuyau mesure environ $\res{447}$ m.}{$1$ hm $=100$ m}
\cas{4)}
\begin{calculs}
\vect{MR} & \coord{-6}{2} &&\\
MR &= \sqrt{36+4} &&\\
MR &= 2\sqrt{10} &&\\
MR &\approx 6{,}325 &&
\end{calculs}
\explique{$MR$ vaut environ $\res{632}$ m, plus que $MH$ : le projeté orthogonal donne bien \res{le chemin le plus court}.}{$R$ est sur $(d)$ : $0-8+8=0$}
\end{exo}

\partie{Calcul d'un angle}

\begin{exo}[Calculer un angle]
\cas{a)}
\begin{calculs}
\vec{u}\cdot\vec{v} &= 2\times 1+0\times\sqrt{3} &&\justif{$xx'+yy'$}\\
\vec{u}\cdot\vec{v} &= 2 &&\\
\norme{\vec{u}} &= 2 &&\justif{$\sqrt{4+0}$}\\
\norme{\vec{v}} &= 2 &&\justif{$\sqrt{1+3}$}\\
\cos(\alpha) &= \frac{2}{2\times 2} &&\justif{propriété 13}\\
\cos(\alpha) &= \frac{1}{2} &&
\end{calculs}
\explique{L'angle mesure $\res{\alpha=60^{\circ}}$.}{valeur remarquable}
\cas{b)}
\begin{calculs}
\vect{AB} & \coord{5}{0} &&\justif{les deux vecteurs issus de $A$}\\
\vect{AC} & \coord{2}{3} &&\\
\vect{AB}\cdot\vect{AC} &= 5\times 2+0\times 3 &&\\
\vect{AB}\cdot\vect{AC} &= 10 &&\\
AB &= 5 &&\\
AC &= \sqrt{13} &&\justif{$\sqrt{4+9}$}\\
\cos\big(\widehat{BAC}\big) &= \frac{10}{5\sqrt{13}} &&\justif{propriété 13}\\
\cos\big(\widehat{BAC}\big) &= \frac{2}{\sqrt{13}} &&\\
\widehat{BAC} &\approx \res{56{,}3^{\circ}} &&\justif{touche $\cos^{-1}$, en mode degré}
\end{calculs}
\cas{c)}
\explique{Pas de coordonnées : on calcule le produit scalaire avec les normes.}{méthode 13, sommet commun $K$}
\begin{calculs}
\vect{KL}\cdot\vect{KM} &= \frac{1}{2}\left(KL^{2}+KM^{2}-LM^{2}\right) &&\justif{propriété 10}\\
\vect{KL}\cdot\vect{KM} &= \frac{1}{2}\left(49+64-169\right) &&\\
\vect{KL}\cdot\vect{KM} &= -28 &&\\
\cos\big(\widehat{LKM}\big) &= \frac{-28}{7\times 8} &&\justif{propriété 13}\\
\cos\big(\widehat{LKM}\big) &= -\frac{1}{2} &&
\end{calculs}
\explique{L'angle mesure $\res{\widehat{LKM}=120^{\circ}}$.}{cosinus négatif : angle obtus}
\end{exo}

\begin{exo}[Les trois angles d'un triangle]
\explique{On calcule d'abord les coordonnées des vecteurs et les longueurs.}{les mêmes servent pour les trois angles}
\begin{calculs}
\vect{AB} & \coord{4}{-2} &&\\
\vect{AC} & \coord{3}{2} &&\\
\vect{BC} & \coord{-1}{4} &&\\
AB &= \sqrt{20} &&\justif{$\sqrt{16+4}$}\\
AC &= \sqrt{13} &&\justif{$\sqrt{9+4}$}\\
BC &= \sqrt{17} &&\justif{$\sqrt{1+16}$}
\end{calculs}
\cas{Angle en $A$}
\begin{calculs}
\vect{AB}\cdot\vect{AC} &= 12-4 &&\justif{$4\times 3+(-2)\times 2$}\\
\vect{AB}\cdot\vect{AC} &= 8 &&\\
\cos\big(\widehat{BAC}\big) &= \frac{8}{\sqrt{20}\times\sqrt{13}} &&\\
\widehat{BAC} &\approx \res{60{,}3^{\circ}} &&\justif{touche $\cos^{-1}$}
\end{calculs}
\cas{Angle en $B$}
\begin{calculs}
\vect{BA}\cdot\vect{BC} &= -4\times(-1)+2\times 4 &&\justif{$\vect{BA}\coord{-4}{2}$, l'opposé de $\vect{AB}$}\\
\vect{BA}\cdot\vect{BC} &= 12 &&\\
\cos\big(\widehat{ABC}\big) &= \frac{12}{\sqrt{20}\times\sqrt{17}} &&\\
\widehat{ABC} &\approx \res{49{,}4^{\circ}} &&
\end{calculs}
\cas{Angle en $C$}
\begin{calculs}
\vect{CA}\cdot\vect{CB} &= -3\times 1+(-2)\times(-4) &&\justif{$\vect{CA}\coord{-3}{-2}$ et $\vect{CB}\coord{1}{-4}$}\\
\vect{CA}\cdot\vect{CB} &= 5 &&\\
\cos\big(\widehat{ACB}\big) &= \frac{5}{\sqrt{13}\times\sqrt{17}} &&\\
\widehat{ACB} &\approx \res{70{,}3^{\circ}} &&
\end{calculs}
\explique{$60{,}3+49{,}4+70{,}3$ vaut $180$ : la somme des angles vaut \res{environ $180^{\circ}$}.}{les arrondis peuvent décaler la somme d'un dixième}
\end{exo}

\begin{exo}[Deux expressions du produit scalaire]
\cas{1)}
\begin{calculs}
\vect{AB} & \coord{6}{0} &&\\
\vect{AC} & \coord{3}{3} &&\\
\vect{AB}\cdot\vect{AC} &= 6\times 3+0\times 3 &&\justif{$xx'+yy'$}\\
\vect{AB}\cdot\vect{AC} &= \res{18} &&
\end{calculs}
\cas{2)}
\begin{calculs}
AB &= \res{6} &&\\
AC &= \sqrt{18} &&\justif{$\sqrt{9+9}$}\\
AC &= \res{3\sqrt{2}} &&\\
\vect{AB}\cdot\vect{AC} &= AB\times AC\times\cos\big(\widehat{BAC}\big) &&\justif{définition 4}\\
\vect{AB}\cdot\vect{AC} &= \res{18\sqrt{2}\cos\big(\widehat{BAC}\big)} &&\justif{$6\times 3\sqrt{2}$}
\end{calculs}
\cas{3)}
\begin{calculs}
18\sqrt{2}\cos\big(\widehat{BAC}\big) &= 18 &&\justif{les deux expressions sont égales}\\
\cos\big(\widehat{BAC}\big) &= \frac{1}{\sqrt{2}} &&\justif{on divise par $18\sqrt{2}$}\\
\cos\big(\widehat{BAC}\big) &= \frac{\sqrt{2}}{2} &&
\end{calculs}
\explique{L'angle mesure $\res{\widehat{BAC}=45^{\circ}}$.}{valeur remarquable}
\cas{4)}
\begin{calculs}
\vect{CA} & \coord{-3}{-3} &&\\
\vect{CB} & \coord{3}{-3} &&\\
\vect{CA}\cdot\vect{CB} &= -9+9 &&\\
\vect{CA}\cdot\vect{CB} &= 0 &&\\
CA &= \sqrt{18} &&\\
CB &= \sqrt{18} &&\justif{$\sqrt{9+9}$}
\end{calculs}
\explique{Le triangle $ABC$ est \res{rectangle en $C$} (produit scalaire nul) et \res{isocèle en $C$} ($CA=CB$).}{on retrouve l'angle de $45^{\circ}$ en $A$}
\end{exo}

\partie{Théorème d'Al-Kashi}

\begin{exo}[Calculer une longueur]
\cas{a)}
\begin{calculs}
BC^{2} &= AB^{2}+AC^{2}-2\times AB\times AC\times\cos\big(\widehat{BAC}\big) &&\justif{on recopie le théorème}\\
BC^{2} &= 49+225-2\times 7\times 15\times\frac{1}{2} &&\justif{$\cos(60^{\circ})=\frac{1}{2}$}\\
BC^{2} &= 169 &&\justif{$274-105$}\\
BC &= \res{13} &&\justif{une longueur est positive}
\end{calculs}
\cas{b)}
\begin{calculs}
FG^{2} &= EF^{2}+EG^{2}-2\times EF\times EG\times\cos\big(\widehat{FEG}\big) &&\justif{angle en $E$}\\
FG^{2} &= 25+9-2\times 5\times 3\times\left(-\frac{1}{2}\right) &&\justif{$\cos(120^{\circ})=-\frac{1}{2}$}\\
FG^{2} &= 49 &&\justif{$34+15$ : moins par moins}\\
FG &= \res{7} &&
\end{calculs}
\cas{c)}
\begin{calculs}
EF^{2} &= DE^{2}+DF^{2}-2\times DE\times DF\times\cos\big(\widehat{EDF}\big) &&\justif{angle en $D$}\\
EF^{2} &= 2+25-2\times\sqrt{2}\times 5\times\frac{\sqrt{2}}{2} &&\justif{$\big(\sqrt{2}\big)^{2}=2$}\\
EF^{2} &= 27-10 &&\justif{$\sqrt{2}\times\sqrt{2}=2$}\\
EF^{2} &= 17 &&\\
EF &= \res{\sqrt{17}} &&
\end{calculs}
\cas{d)}
\begin{calculs}
ST^{2} &= RS^{2}+RT^{2}-2\times RS\times RT\times\cos\big(\widehat{SRT}\big) &&\justif{angle en $R$}\\
ST^{2} &= 3+4-2\times\sqrt{3}\times 2\times\left(-\frac{\sqrt{3}}{2}\right) &&\justif{$\cos(150^{\circ})=-\frac{\sqrt{3}}{2}$}\\
ST^{2} &= 7+6 &&\justif{$2\times\sqrt{3}\times\sqrt{3}=6$}\\
ST^{2} &= 13 &&\\
ST &= \res{\sqrt{13}} &&
\end{calculs}
\cas{e)}
\begin{calculs}
LM^{2} &= KL^{2}+KM^{2}-2\times KL\times KM\times\cos\big(\widehat{LKM}\big) &&\justif{angle en $K$}\\
LM^{2} &= 16+36-48\times\frac{\sqrt{3}}{2} &&\justif{$\cos(30^{\circ})=\frac{\sqrt{3}}{2}$}\\
LM^{2} &= 52-24\sqrt{3} &&\\
LM &= \sqrt{52-24\sqrt{3}} &&\\
LM &\approx \res{3{,}23} &&\justif{calculatrice}
\end{calculs}
\end{exo}

\begin{exo}[Un angle avec Al-Kashi]
\cas{1)}
\begin{calculs}
BC^{2} &= AB^{2}+AC^{2}-2\times AB\times AC\times\cos\big(\widehat{BAC}\big) &&\justif{théorème d'Al-Kashi}\\
2\times AB\times AC\times\cos\big(\widehat{BAC}\big) &= AB^{2}+AC^{2}-BC^{2} &&\justif{on échange les termes de côté}\\
\cos\big(\widehat{BAC}\big) &= \res{\dfrac{AB^{2}+AC^{2}-BC^{2}}{2\times AB\times AC}} &&\justif{on divise par $2\times AB\times AC$}
\end{calculs}
\cas{2)}
\begin{calculs}
\cos\big(\widehat{BAC}\big) &= \dfrac{25+36-49}{2\times 5\times 6} &&\\
\cos\big(\widehat{BAC}\big) &= \dfrac{12}{60} &&\\
\cos\big(\widehat{BAC}\big) &= \dfrac{1}{5} &&\\
\widehat{BAC} &\approx \res{78{,}5^{\circ}} &&\justif{touche $\cos^{-1}$, en mode degré}
\end{calculs}
\end{exo}

\begin{exo}[Largeur d'un lac]
\cas{1)}
\begin{calculs}
AB^{2} &= SA^{2}+SB^{2}-2\times SA\times SB\times\cos\big(\widehat{ASB}\big) &&\justif{théorème d'Al-Kashi, angle en $S$}\\
AB^{2} &= 90\,000+250\,000-2\times 300\times 500\times\frac{1}{2} &&\\
AB^{2} &= 190\,000 &&\justif{$340\,000-150\,000$}\\
AB &= \sqrt{190\,000} &&\\
AB &= \res{100\sqrt{19}} &&\justif{$\sqrt{10\,000}\times\sqrt{19}$}
\end{calculs}
\explique{La distance $AB$ vaut environ $\res{436}$ m.}{$100\sqrt{19}\approx 435{,}9$}
\cas{2)}
\explique{Le théorème de Pythagore ne vaut que dans un triangle \res{rectangle} ; ici l'angle en $S$ mesure $60^{\circ}$.}{Al-Kashi généralise Pythagore : avec un angle droit, $\cos(90^{\circ})=0$ et le dernier terme disparaît}
\end{exo}

\begin{exo}[Diagonales et angle]
\cas{1)}
\begin{calculs}
\vect{AB}\cdot\vect{AD} &= 6\times 4\times\cos(60^{\circ}) &&\justif{définition 4}\\
\vect{AB}\cdot\vect{AD} &= \res{12} &&
\end{calculs}
\cas{2)}
\begin{calculs}
BD^{2} &= AB^{2}+AD^{2}-2\times AB\times AD\times\cos\big(\widehat{DAB}\big) &&\justif{triangle $ABD$, angle en $A$}\\
BD^{2} &= 36+16-2\times 12 &&\justif{$AB\times AD\times\cos(60^{\circ})=12$}\\
BD^{2} &= 28 &&\\
BD &= \res{2\sqrt{7}} &&\justif{$\sqrt{4}\times\sqrt{7}$}
\end{calculs}
\cas{3)}
\explique{Dans un parallélogramme, deux angles consécutifs ont pour somme $180^{\circ}$ : $\widehat{ABC}=180^{\circ}-60^{\circ}$, soit $\res{120^{\circ}}$.}{$(AD)$ et $(BC)$ sont parallèles}
\begin{calculs}
AC^{2} &= AB^{2}+BC^{2}-2\times AB\times BC\times\cos\big(\widehat{ABC}\big) &&\justif{triangle $ABC$, angle en $B$, $BC=AD=4$}\\
AC^{2} &= 36+16-2\times 6\times 4\times\left(-\frac{1}{2}\right) &&\\
AC^{2} &= 76 &&\justif{$52+24$}\\
AC &= \res{2\sqrt{19}} &&\justif{$\sqrt{4}\times\sqrt{19}$}
\end{calculs}
\cas{4)}
\begin{calculs}
AC^{2}+BD^{2} &= 76+28 &&\\
AC^{2}+BD^{2} &= 104 &&\\
2\left(AB^{2}+AD^{2}\right) &= 2\times 52 &&\\
2\left(AB^{2}+AD^{2}\right) &= 104 &&
\end{calculs}
\explique{\res{L'égalité est vérifiée.}}{c'est la même relation qu'à l'exercice 32}
\end{exo}

\partie{Ensemble des points \texorpdfstring{$M$ tels que $\vect{MA}\cdot\vect{MB}=0$}{M tels que MA.MB = 0}}

\begin{exo}[Cercle de diamètre donné]
\cas{a)}
\explique{$(\mathcal{C})$ est le cercle de diamètre $[AB]$ : son centre est le milieu $I$ de $[AB]$, son rayon est $IA$.}{théorème 7}
\begin{calculs}
I & \left(\dfrac{-3+5}{2}\,;\dfrac{1-3}{2}\right) &&\justif{coordonnées du milieu}\\
I & (1\,;-1) &&\\
IA^{2} &= (-3-1)^{2}+(1+1)^{2} &&\justif{le carré du rayon suffit}\\
IA^{2} &= 20 &&\justif{$16+4$}
\end{calculs}
\explique{Une équation de $(\mathcal{C})$ est $\res{(x-1)^{2}+(y+1)^{2}=20}$.}{théorème 2}
\cas{b)}
\begin{calculs}
(3-1)^{2}+(3+1)^{2} &= 4+16 &&\justif{on remplace $x$ et $y$ par $3$}\\
(3-1)^{2}+(3+1)^{2} &= 20 &&
\end{calculs}
\explique{$C$ \res{est sur $(\mathcal{C})$}, donc $\vect{CA}\cdot\vect{CB}=0$ et le triangle $ABC$ \res{est rectangle en $C$}.}{théorème 7}
\cas{c)}
\explique{$\vect{DM}$ est orthogonal à $\vect{DE}$ : $(\mathcal{D})$ est \res{la droite qui passe par $D$, perpendiculaire à $(DE)$}, de vecteur normal $\vect{DE}$.}{théorème 5}
\begin{calculs}
\vect{DE} & \coord{3}{2} &&\justif{$5-2$ et $1-(-1)$}\\
\vect{DM} & \coord{x-2}{y+1} &&\justif{avec $M(x\,;y)$}\\
\vect{DE}\cdot\vect{DM} &= 3(x-2)+2(y+1) &&\justif{$xx'+yy'$}\\
\vect{DE}\cdot\vect{DM} &= 3x+2y-4 &&\justif{$3x-6+2y+2$}
\end{calculs}
\explique{Une équation de $(\mathcal{D})$ est $\res{3x+2y-4=0}$.}{}
\end{exo}

\begin{exo}[Retrouver le cercle par le calcul]
\cas{1)}
\begin{calculs}
\vect{MA} & \res{\coord{-2-x}{-y}} &&\justif{arrivée $A$, départ $M$}\\
\vect{MB} & \res{\coord{4-x}{-y}} &&
\end{calculs}
\cas{2)}
\begin{calculs}
\vect{MA}\cdot\vect{MB} &= (-2-x)(4-x)+(-y)\times(-y) &&\justif{$xx'+yy'$}\\
\vect{MA}\cdot\vect{MB} &= -8+2x-4x+x^{2}+y^{2} &&\justif{on développe}\\
\vect{MA}\cdot\vect{MB} &= \res{x^{2}-2x+y^{2}-8} &&
\end{calculs}
\cas{3)}
\begin{calculs}
x^{2}-2x+y^{2}-8 &= 0 &&\justif{question 2}\\
(x-1)^{2}-1+y^{2}-8 &= 0 &&\justif{$x^{2}-2x=(x-1)^{2}-1$}\\
(x-1)^{2}+y^{2} &= 9 &&
\end{calculs}
\explique{Donc $\vect{MA}\cdot\vect{MB}=0$ équivaut à $\res{(x-1)^{2}+y^{2}=9}$.}{chaque ligne équivaut à la précédente}
\cas{4)}
\explique{C'est \res{le cercle de centre $(1\,;0)$ et de rayon $3$}.}{théorème 2}
\explique{C'était prévisible : $(1\,;0)$ est le milieu de $[AB]$ et $AB=6$ ; c'est \res{le cercle de diamètre $[AB]$}.}{théorème 7}
\end{exo}

\begin{exo}[Angle droit sur l'axe des abscisses]
\cas{1)}
\begin{calculs}
\vect{MA} & \coord{1-x}{3} &&\justif{$3-0$}\\
\vect{MB} & \coord{5-x}{1} &&\\
\vect{MA}\cdot\vect{MB} &= (1-x)(5-x)+3\times 1 &&\\
\vect{MA}\cdot\vect{MB} &= 5-x-5x+x^{2}+3 &&\justif{on développe}\\
\vect{MA}\cdot\vect{MB} &= \res{x^{2}-6x+8} &&
\end{calculs}
\cas{2)}
\begin{calculs}
\Delta &= 36-32 &&\justif{$(-6)^{2}-4\times 1\times 8$}\\
\Delta &= 4 &&\\
x_{1} &= \dfrac{6-2}{2} &&\\
x_{1} &= 2 &&\\
x_{2} &= \dfrac{6+2}{2} &&\\
x_{2} &= 4 &&
\end{calculs}
\explique{Le triangle $AMB$ est rectangle en $M$ pour $\res{M_{1}(2\,;0)}$ et $\res{M_{2}(4\,;0)}$.}{rectangle en $M$ : $\vect{MA}\cdot\vect{MB}=0$}
\cas{3)}
\begin{calculs}
I & \left(\dfrac{1+5}{2}\,;\dfrac{3+1}{2}\right) &&\justif{centre : milieu de $[AB]$}\\
I & (3\,;2) &&\\
IA^{2} &= (1-3)^{2}+(3-2)^{2} &&\\
IA^{2} &= 5 &&
\end{calculs}
\explique{Le cercle de diamètre $[AB]$ a pour équation $\res{(x-3)^{2}+(y-2)^{2}=5}$.}{}
\begin{calculs}
(2-3)^{2}+(0-2)^{2} &= 1+4 &&\justif{point $M_{1}$}\\
(2-3)^{2}+(0-2)^{2} &= 5 &&\\
(4-3)^{2}+(0-2)^{2} &= 1+4 &&\justif{point $M_{2}$}\\
(4-3)^{2}+(0-2)^{2} &= 5 &&
\end{calculs}
\explique{Les deux points \res{sont sur ce cercle}.}{ce sont les points où le cercle coupe l'axe des abscisses}
\cas{4)}
\begin{calculs}
\vect{AB} & \coord{4}{-2} &&\\
\vect{AN} & \coord{x-1}{y-3} &&\justif{avec $N(x\,;y)$}\\
\vect{AB}\cdot\vect{AN} &= 4(x-1)-2(y-3) &&\\
\vect{AB}\cdot\vect{AN} &= 4x-2y+2 &&\justif{$4x-4-2y+6$}
\end{calculs}
\explique{L'ensemble est la droite qui passe par $A$, perpendiculaire à $(AB)$, d'équation $\res{4x-2y+2=0}$, ou $2x-y+1=0$.}{on divise par $2$}
\end{exo}

\end{document}
